\documentclass[reqno]{amsart}
\usepackage{a4wide}
\usepackage{hyperref}

\AtBeginDocument{{\noindent\small
\emph{Electronic Journal of Differential Equations},
Vol. 2026 (2026), No. 03, pp. 1--27.\newline
ISSN: 1072-6691. URL: https://ejde.math.txstate.edu, DOI: 10.58997/ejde.2026.03}
\thanks{\copyright 2026. This work is licensed under a CC BY 4.0 license.}
\vspace{8mm}}

\begin{document}
\title[\hfilneg EJDE-2026/03\hfil 
Fractional Kirchhoff-Schr\"odinger-Poisson system]
{Normalized solutions for a fractional Kirchhoff-Schr\"odinger-Poisson
systems with \\ critical growth}

\author[Y.-Q. Zhao, J.-F. Liao \hfil EJDE-2026/03\hfilneg]
{Yu-Qin Zhao, Jia-Feng Liao}

\address{Yu-Qin Zhao \newline
School of Mathematics Sciences,
China West Normal University, 
Nanchong, Sichuan 637009, China}
\email{2943658736@qq.com}

\address{Jia-Feng Liao (corresponding author) \newline
School of Mathematics Sciences,
China West Normal University, 
Nanchong, Sichuan 637009, China}
\email{liaojiafeng@163.com}

\thanks{Submitted September 18, 2025. Published January 10, 2026.}
\subjclass[2020]{35B33, 35R11, 35J50}
\keywords{Kirchhoff-Schr\"odinger-Poisson system;  normalized solution; 
variational methods; 
\hfill\break\indent  $L^2$-subcritical; $L^2$-supercritical}

\begin{abstract}
In this article, we study the fractional Kirchhoff-Schr\"odinger-Poisson
system with Sobolev critical growth
\begin{gather*}
\Big(a+b\int_{\mathbb{R}^3}|(-\Delta)^{s/2}u|^2dx\Big)
(-\Delta)^s u+ \phi u=\lambda u+ \mu|u|^{p-2}u+|u|^{2^*_s-2}u,
\quad \text{in }\mathbb{R}^3, \\
(-\Delta)^s \phi=u^2, \quad  \text{in }\mathbb{R}^3,
\end{gather*}
where $a,b>0$, $s\in(\frac{3}{4},1)$,   $p \in(2,2^*_s)$, and $\mu>0$ is a parameter,
$\lambda \in \mathbb{R}$ is an undermined parameter. For this problem, under the
$L^2$-subcritical, $p\in(2, \frac{4s}{3}+2)$, we obtain the multiplicity of the
normalized solutions by means of the truncation technique, concentration-compactness
principle, and  genus theory.
In the $L^2$-supercritical, $p\in( \frac{8s}{3}+2, 2^*_s)$, we prove a couple
of normalized solutions by developing a fiber map and using the
concentration-compactness principle.
\end{abstract}

\maketitle
\numberwithin{equation}{section}
\newtheorem{theorem}{Theorem}[section]
\newtheorem{lemma}[theorem]{Lemma}
\newtheorem{corollary}[theorem]{Corollary}
\newtheorem{proposition}[theorem]{Proposition}
\newtheorem{remark}[theorem]{Remark}
\allowdisplaybreaks

\section{Introduction}

In this article, we study the nonlinear Kirchhoff-Schr\"odinger-Poisson system
with Sobolev critical growth
\begin{equation}\label{lab1.1}
\begin{gathered}
(a+b\int_{\mathbb{R}^3}|(-\Delta)^{s/2}u|^2dx)(-\Delta)^s u+ \phi u
= \lambda u+\mu|u|^{p-2}u+|u|^{{2^*_s}-2}u,  \quad \text{in }\mathbb{R}^3, \\
(-\Delta)^s \phi=u^2, \quad   \text{in }\mathbb{R}^3,
\end{gathered}
\end{equation}
where $a,b>0$, $s\in(\frac{3}{4},1)$, and  $p \in(2,2^*_s)$,  $\mu>0$ and
$\lambda \in \mathbb{R}$ are parameters.
Here $(-\Delta)^s$ $(s\in(0,1))$ is the fractional Laplacian operator which is
defined by
$$
(-\Delta)^s u(x)=C_s {\rm P.V.}\int_{\mathbb{R}^3}\frac{u(x)-u(y)}{|x-y|^{3+2s}}dy
=C_s\lim_{\epsilon\to 0}\int_{{\mathbb{R}^3}
\setminus B_\epsilon(x)}\frac{u(x)-u(y)}{|x-y|^{3+2s}}dy,
$$
for $u\in \mathcal{S}(\mathbb{R}^3)$, where $\mathcal{S}(\mathbb{R}^3) $ is the
Schwartz space  of rapidly decaying $C^\infty$ functions, $B_\epsilon(x)$ denote an
open ball of radius $\epsilon$ at $x$ and $C_s$ is a normalization constant.

System \eqref{lab1.1} has been motivated by the
time-dependent fractional Schr\"odinger-Poisson system
\begin{equation}\label{lab1.12}
\begin{gathered}
i\frac{\partial \Psi}{\partial \tau } =(-\Delta)^s\Psi +\lambda\phi\Psi-f(x,|\Psi|),
\quad  x\in \mathbb{R}^3,\\
(-\Delta)^t\phi=|\Psi|^2, \quad  x\in \mathbb{R}^3,
\end{gathered}
\end{equation}
where $\Psi: {\mathbb{R}\times \mathbb{R}^3}\to \mathbb{C}$, $s,t \in(0,1)$,
$\lambda \in \mathbb{R}$. It is well-known that, the first equation in
system \eqref{lab1.12} was used by Laskin (see \cite{L2998, L3998}) to extend the Feynman path integral, from Brownian-like to L$\mathrm {\acute{e}}$vy-like quantum mechanical paths. This class of fractional Schr\"odinger equations with a repulsive nonlocal Coulombic potential can be approximated by the Hartree-Fock equations to describe a quantum mechanical system of many particles. For more application backgrounds on the fractional Laplacian see \cite{C4998, D7998, M6998}.

When looking for solutions to system (1.1), there are two distinct options  regarding  the frequency parameter  $\lambda$. One is to regard the frequency $\lambda$ as a given constant. Xiang and Wang \cite{H2998} first investigated  the  fractional  Kirchhoff-Schr\"odinger-Poisson system, and obtained the existence, multiplicity and asymptotic  behavior  of  nonnegative  solutions. In recent years, researchers have shown growing interest in the fractional Kirchhoff-Schr\"odinger-Poisson system
\begin{equation}\label{lab1.156}
\begin{gathered}
(a+b\int_{\mathbb{R}^3}|(-\Delta)^{s/2}u|^2dx)(-\Delta)^s u+V(x)u+\mu \phi u= f(x,u),   
\quad  \text{in }\mathbb{R}^3, \\
(-\Delta)^t \phi=\mu u^2, \quad   \text{in }\mathbb{R}^3,
\end{gathered}
\end{equation}
where $a>0$, $b\geq 0$. When $V(x)=0$, $f(x,u)=f(u)$, by utilizing minimax argument, 
Ambrosio \cite{A1111} obtained the existence of a nontrivial solutions  
for system \eqref{lab1.156} with Berestycki-Lions type nonlinearities. 
When   $V(x)\neq 0$, Wang et.al \cite{W1112} studied the existence of ground solutions 
for system \eqref{lab1.156} with  $V(x)=1$ and 
$f(x,u)=(|x|^{-\theta}\ast F(u))f(u)$, $\theta \in(0, 3-2t)$, and used the Pohoz\v{a}ev 
type  manifold; Then, under some assumptions on $V$ and $f$, by using constraint 
variational approach and a quantitative deformation lemma, Meng et.al \cite{M1113} 
proved the existence of  the least energy sign-changing solutions for system 
\eqref{lab1.156}.  Feng et al.\ \cite{F1114} applied a similar method studied the 
least energy sign-changing solutions to the following critical fractional 
Kirchhoff-Schr\"odinger-Poisson system with steep potential well
  %\label{lab1.1156}
\begin{gather*}
\Big(a+b\int_{\mathbb{R}^3}|(-\Delta)^{s/2} u|^2dx\Big)
(-\Delta)^s u+V_{\lambda}(x)u+ \phi u=|u|^{p-2}u+|u|^{2^*_s-2}u, \quad
 \text{in }\mathbb{R}^3, \\
(-\Delta)^t \phi=u^2, \quad   \text{in }\mathbb{R}^3,
\end{gather*}
where $s\in(\frac{3}{4}, 1)$, $t\in(0,1)$, $V_\lambda(x)=\lambda V(x)+1$,
$\lambda>0$ and $p>4$.  Jian et al.\ \cite{J1115}, deal with the
fractional Kirchhoff-Schr\"odinger-Poisson system with steep potential well
 % \label{lab1.1567}
\begin{gather*}
(a+b\int_{\mathbb{R}^3}|(-\Delta)^{s/2}u|^2dx)(-\Delta)^s u+\lambda V(x)u
 +\mu \phi u= |u|^{p-2}u,  \quad \text{in }\mathbb{R}^3, \\
(-\Delta)^t \phi= u^2, \quad   \text{in }\mathbb{R}^3,
\end{gather*}
where $s\in[\frac{3}{4}, 1)$, $t\in(0,1)$, $2<p<4$, $a>0$ is a constant, and
$b, \lambda, \mu$ are positive parameters. By applying the truncation technique and
the parameter-dependent compactness lemma, they first proved the existence of
positive solutions.  Furthermore, they investigated the   asymptotic behavior as
$b\to 0$, $\lambda \to \infty$ and $\mu \to 0$, respectively.
For other existence results, we refer to \cite{F1125, W1145, W1135}  and the
references therein.

When $a=1$, $b=0$, system \eqref{lab1.156} reduces to the following fractional
Schr\"odinger-Poisson system
\begin{equation}\label{lab1.15776}
\begin{gathered}
(-\Delta)^s u+V(x)u+\mu \phi u= f(x,u),   \quad \text{in }\mathbb{R}^3, \\
(-\Delta)^t \phi=\mu u^2, \quad   \text{in }\mathbb{R}^3.
\end{gathered}
\end{equation}
Recently, under various potentials and nonlinear terms, most scholars have
investigated the existence and multiplicity of ground state  solutions,
sign-changing solutions, and nontrivial solutions for system \eqref{lab1.15776}.
 For further details, we refer  the interesting readers to see
\cite{L1159, D1399, G1312, G1314, P1199} and so on.

Alternatively, the other one is to regard the frequency $\lambda$ as an unknown
quantity. In such point of view, it is natural to prescribe the mass, i.e.,
the $L^2$-norm, so that $\lambda$ can be interpreted as a Lagrange multiplier.
Solutions of this type are often referred to as normalized solutions.
Nowadays, from a physical point of view, some physicists are very interested in the
 normalized solutions, see for example \cite{Bartsch99, luo288,  yang1288}.
There a few results are related to the study of normalized solutions for the
fractional Kirchhoff-Schr\"odinger-Poisson system except Wang et al.\ \cite{wang1388}.
They considered the  fractional Kirchhoff-Schr\"odinger-Poisson system
\[  %\label{lab1.1234}
\begin{gathered}
\Big(a+b\int_{\mathbb{R}^3}|(-\Delta)^{s/2}u|^2dx\Big)(-\Delta)^s u+ \phi u= f
(u)+\lambda u    \quad \text{in }\mathbb{R}^3, \\
(-\Delta)^t \phi=u^2 \quad    \text{in }\mathbb{R}^3,
\end{gathered}
\]
where $a,b>0$, $s, t\in (0,1)$, $2s+2t>3$, $\lambda\in \mathbb{R}$,
$f\in C(\mathbb{R}, \mathbb{R})$ satisfies some general conditions which contain
the case  $f(u)\sim|u|^{p-2}u$ with
$p\in(\frac{8s+4t-3}{s+t},\frac{8s}{3}+2)\cup(\frac{8s}{3}+2, 2^*_s)$,
$2^*_s=\frac{6}{3-2s}$. They obtained
the existence of normalized solutions by using the Pohoz\v{a}ev manifold and
variational method.

 Recently, He and Meng \cite{H2588} studied the existence and multiplicity of the
normalized solutions for the nonlinear fractional Schr\"odinger-Poisson system with
Sobolev critical exponent
\[ % \label{lab1.1434}
\begin{gathered}
(-\Delta)^s u+ \alpha\phi u= \mu|u|^{p-2}u+|u|^{2^*_s-2}u+\lambda u   \quad
 \text{in }\mathbb{R}^3, \\
(-\Delta)^t \phi=u^2 \quad    \text{in }\mathbb{R}^3,
\end{gathered}
\]
 where $s, t\in (0,1)$, $2s+2t>3$, $p\in(2, 2^*_s)$, $\alpha, \mu>0$ are parameters
and $\lambda\in \mathbb{R}$.

Motivated by \cite{H2588} and \cite{wang1388}, a natural question is whether the
fractional Kirchhoff-Schr\"odinger-Poisson system with the Sobolev critical growth can
be applied to obtain the  existence and multiplicity   of normalized solutions for $p$
in distinct ranges.
Therefore, we study system \eqref{lab1.1} and give an affirmative answer.
In addition, to recover the compactness, we will take $H^s_r(\mathbb{R}^3)$ as a
working space.

By using the Lax-Milgram theorem, for any $u\in H^s_r(\mathbb{R}^3)$, a unique
$\phi ^s_u(x)\in D^{s,2}(\mathbb{R}^3)$ is given by
\[ \label{lab1.4}
\phi ^s_u(x)=|x|^{2s-3}\ast|u|^2=\int_{\mathbb{R}^3}\frac{u^2(y)}{|x-y|^{3-2s}}dy,
\]
such that $(-\Delta)^s\phi =u^2$ and that inserting it into the first equation of
system \eqref{lab1.1}, then system \eqref{lab1.1} can be transformed into the following
single equation
\begin{equation}\label{lab1.5}
\Big(a+b\int_{\mathbb{R}^3}|(-\Delta)^{s/2}u|^2dx\Big)(-\Delta)^s u
 + \phi^s_u u= \lambda u+ \mu|u|^{p-2}u+ |u|^{2^*_s-2}u, \quad\text{in }\mathbb{R}^3.
\end{equation}
It can be proved  that to find the normalized solutions of system \eqref{lab1.1}
is to seek the critical points of the functional
\begin{align*}
  I_\mu(u)
&=\frac{a}{2}\int_{\mathbb{R}^3}|(-\Delta)^{s/2}u|^2dx
  +\frac{b}{4}\Big(\int_{\mathbb{R}^3}|(-\Delta)^{s/2}u|^2dx\Big)^2\\
&\quad +\frac{1}{4}\int_{\mathbb{R}^3}\phi ^s_uu^2dx
 -\frac{\mu}{p}\int _{\mathbb{R}^3}|u|^pdx
 -\frac{1}{2^*_s}\int _{\mathbb{R}^3}|u|^{2^*_s}dx,
\end{align*}
under the constraint
\[ \label{1.8}
S_r(c)=\Big\{u \in H^s_r(\mathbb{R}^3):\int_{\mathbb{R}^3}|u|^2dx=c^2,c>0\Big\}.
\]
It is well known that $I_\mu\in {C}^1(H^s_r(\mathbb{R}^3), \mathbb{R})$.
Here are our main results.

\begin{theorem}\label{thm1.1}
Assume $a,b>0$, $s\in(\frac{3}{4},1)$ and $2<p<\frac{4s}{3}+2$, for given
$k\in \mathbb{N}$, then there exist $\beta>0$,
$\Lambda=\big(\frac{1}{2}-\frac{1}{2^*_s}\big) a^{\frac{3}{2s}}S^{\frac{3}{2s}}>0$,
$D=\big(\frac{1}{p}-\frac{1}{2^*_s}\big)C(p,s)R_1^{p\delta_{p,s}}>0$ independent of
$k$ and $\mu^*_k>0$ large, such that system \eqref{lab1.1}  possesses at least $k$
couples $(u_j, \lambda_j)\in H^s_r(\mathbb{R}^3)\times \mathbb{R}$ of weak solutions
for $\mu>\mu^*_k$ and
\[  %\label{2.4}
c\in\Big(0, \min\Big\{\Big(\frac{\beta}{\mu}\Big)^{\frac{1}{p(1-\delta_{p,s})}},
\Big(\frac{\Lambda}{D\mu}\Big)^{\frac{1}{p(1-\delta_{p,s})}}\Big\}\Big)
\]
with $\int_{\mathbb{R}^3}|u_j|^2dx=c^2$, $\lambda_j<0$ for all $j=1,2,...,k$,
$\delta_{p,s}=\frac{3(p-2)}{2ps}$.
\end{theorem}

\begin{theorem}\label{thm1.2}
Assume $a,b>0$, $s\in(\frac{3}{4},1)$ and $2+\frac{8s}{3}<p<2^*_s$, then there exists
$\mu^*=\mu^*(c)>0$ large, such that as $\mu>\mu^*$, system \eqref{lab1.1}
possesses a couple $(u_c, \lambda)\in H^s_r(\mathbb{R}^3)\times \mathbb{R}$ of weak
solutions with $\int_{\mathbb{R}^3}|u_c|^2dx=c^2$, $\lambda<0$.
\end{theorem}

\begin{remark}\rm
To our best knowledge, our results are up to date. On the one hand,
the Sobolev critical exponent leads to the lack of compactness.
Even the embedding of the radially symmetric space of $H^s_r(\mathbb{R}^3)$ into
$L^{2^*_s}(\mathbb{R}^3)$ is not compact. Furthermore,
$H^s_r(\mathbb{R}^3)\hookrightarrow L^2(\mathbb{R}^3)$ is also not compact.
Then, the weak limit of Palais-Smale sequences could leave the constrained manifold
$S_r(c)$. Hence, we need to estimate finely the Lagrange multipliers, which is vital
in obtaining compactness. We shall employ the concentration-compactness principle,
mountain pass theorem and the truncation method to overcome the loss of compactness
caused by the critical growth.

On the other hand, no matter $2<p<{\frac{4s}{3}+2}$ or $\frac{8s}{3}+2<p<2^*_s$,
$I_\mu$ on the constrained manifold $S_r(c)$ is all unbounded from below.
Hence, it is unlikely to obtain a solution to system \eqref{lab1.1} by minimizing
method. We adopt some ideas from \cite{A5878} to overcome the difficulty.
\end{remark}

This article is structured as follows:
in section 2, we presents some preliminary results that will be  used frequently
 in the sequel.
Theorem \ref{thm1.1} is proved in section 3, which presents the multiplicity of
normalized solutions for system (1.1) when $p\in(2, \frac{4s}{3}+2)$.
In this section, we address three main challenges. First, we employ the truncation
technique to establish the boundedness of the (PS) sequence. Subsequently, we apply
the concentration compactness principle to restore the compactness lost of
the (PS) sequence due to the critical growth. Finally, we use genus theory to prove
the multiplicity of normalized solutions for system (1.1).
In section 4, when the parameter $\mu>0$ is large, we give another existence result
for system (1.1) with $p\in( \frac{8s}{3}+2, 2^*_s)$ by using the fiber map and
concentration-compactness principle.

\noindent\textbf{Notation}
Throughout this paper, we denote $\|\cdot\|_q$ the usual norm of the  space
$L^q(\mathbb{R}^3)$, $1\leq q <\infty$, $B_r(x)$ denotes the open ball with center at
$x$ and radius $r$, $C$ or $C_i (i=1,2,...)$ denote various positive constants whose
exact values are irrelevant. $\rightharpoonup$ and $\to$ mean the weak and strong
convergence.

\section{Preliminaries}

In this section, we first introduce  some  notations. For any $s\in(0,1)$, the
homogeneous Sobolev space $D^{s,2}(\mathbb{R}^3)$ is defined by
$D^{s,2}(\mathbb{R}^3)=\{u\in L^{2^*_s}(\mathbb{R}^3):\|u\|_{D^{s,2}}<\infty\}$,
where
$$
\|u\|^2_{D^{s,2}}=\int_{\mathbb{R}^3}|(-\Delta)^{s/2}u|^2dx
=\int_{\mathbb{R}^3}\int_{\mathbb{R}^3}\frac{|u(x)-u(y)|^2}{|x-y|^{3+2s}}\,dx\,dy.
$$
The fractional space $H^s(\mathbb{R}^3)$ is defined by
$$
H^s(\mathbb{R}^3)=\{u\in L^2(\mathbb{R}^3): \|u\|_{D^{s,2}}< \infty \},
$$
endowed with the norm
$$
\|u\|^2_{H^s}=\|u\|^2_2+\|u\|^2_{D^{s,2}}.
$$
The best fractional Sobolev constant $S$ is defined as
\begin{equation}\label{2.111}
S=\inf_{u\in D^{s,2}, u\neq 0}
\frac{\|(-\Delta)^{s/2}u\|^2_2}
{(\int_{\mathbb{R}^3}|u|^{2^*_s}dx)^{2/2^*_s}}.
\end{equation}
The work space is
 $$
 H^s_r(\mathbb{R}^3)=\{u \in H^s(\mathbb{R}^3)|~ u(x)=u(|x|)\}.
 $$
Let $\mathbb{H}=H^s_r(\mathbb{R}^3) \times \mathbb{R}$ with the scalar product
$\langle\cdot, \cdot\rangle_{H^s_r}+\langle\cdot, \cdot\rangle_{\mathbb{R}}$, and
the corresponding norm
$\|(\cdot, \cdot)\|^2_\mathbb{H}=\|\cdot\|^2_{H^s_r}+|\cdot|^2_{\mathbb{R}}$.

\begin{proposition}[Hardy-Littlewood-Sobolev inequality] \label{pro2.1}
Let $r,l>1$ and $0<\lambda <N$ with $\frac{1}{r}+\frac{1}{l}+\frac{\lambda}{N}=2$.
Let $f\in L^r(\mathbb{R}^N)$ and $g\in L^l(\mathbb{R}^N)$.
Then there exists a sharp constant $C(N, \lambda, l, r)>0$, independent of $f$ and $g$,
such that
\[
\Big|\int_{\mathbb{R}^N}\int_{\mathbb{R}^N}f(x)g(y)|x-y|^{-\lambda}\,dx\,dy\Big|
\leq C(N, \lambda, l, r)\|f\|_{r}\|g\|_{l}.
\]
Moreover, if $r=l=\frac{2N}{2N-\lambda}$, then
\[
C(N, \lambda, l, r)
=\pi^{\lambda/2}\frac{\Gamma(\frac{N}{2}-\frac{\lambda}{2})}
{\Gamma(N-\frac{\lambda}{2})}\Big(\frac{\Gamma(N)}{\Gamma(\frac{N}{2})}
\Big)^{1-\frac{\lambda}{N}}.
\]
\end{proposition}

From Proposition \ref{pro2.1},  with  $r=l=\frac{6}{3+2s}$, the
 Hardy-Littlewood-Sobolev inequality implies that
\begin{equation}\label{3.3}
\int_{\mathbb{R}^3}\phi^s_u u^2dx=\int_{\mathbb{R}^3}
\left(|x|^{2s-3}\ast u^2\right)u^2dx\leq \Gamma_s\|u\|^4_{\frac{12}{3+2s}},
\end{equation}
where
\[
\Gamma_s=C\big(3, 3-2s, \frac{6}{3+2s}, \frac{6}{3+2s}\big)=\pi^{\frac{3-2s}{2}}
\frac{\Gamma(s)}{\Gamma(\frac{3}{2}+s)}\Big(\frac{\Gamma(3)}
{\Gamma(\frac{3}{2})}\Big)^{\frac{2s}{3}},
\]
where $\Gamma(t)$ is  the Gamma function with $t>0$.
Now, we introduce the Pohoz{\v{a}}ev manifold associated to equation  \eqref{lab1.5},
which can be derived from \cite{T7899}.

\begin{proposition}\label{pro2.2}
Let $u\in H^s_r(\mathbb{R}^3)$ be a weak solution of equation  \eqref{lab1.5},
then $u$ satisfies
\begin{align*}
&\frac{3-2s}{2}\Big(a+b\int_{\mathbb{R}^3}|(-\Delta)^{s/2}u|^2dx\Big)
 \int_{\mathbb{R}^3}|(-\Delta)^{s/2}u|^2dx
 +\frac{3+2s}{4}\int_{\mathbb{R}^3}\phi^s_u u^2dx\\
&=\frac{3}{2}\lambda \int_{\mathbb{R}^3}|u|^2dx
 +\frac{3}{p}\mu \int_{\mathbb{R}^3}|u|^pdx
 +\frac{3}{2^*_s}\int_{\mathbb{R}^3}|u|^{2^*_s}dx.
\end{align*}
\end{proposition}

\begin{lemma}[\cite{D1688}] \label{lem2.1311}
The embedding $H^s_r(\mathbb{R}^3)\hookrightarrow L^q(\mathbb{R}^3)$ is compact for
 any $2< q < 2^*_s$.
\end{lemma}

\begin{lemma}[Fractional Gagliardo-Nirenberg inequality] \label{lem3.2}
Let $0<s<1$, and $p\in(2, 2^*_s)$. Then there exists a constant
$C(p,s)=S^{-\frac{\delta_{p,s}}{2}}>0$ such that
\begin{equation}\label{3.2}
\|u\|^p_p\leq C(p,s)\|(-\Delta)^{s/2}u\|^{\frac{3}{s}(\frac{p}{2}-1)}_2
\|u\|^{\frac{3}{s}(1-\frac{3-2s}{6}p)}_2, \quad \forall u\in H^s(\mathbb{R}^3),
\end{equation}
where $\delta_{p,s}=\frac{3(p-2)}{2ps}$.
\end{lemma}

\begin{lemma}[\cite{D1678}] \label{lem2.171}
If $u_n\rightharpoonup u$ in $H^s_r(\mathbb{R}^3)$, then
\[
\int_{\mathbb{R}^3}\phi^s_{u_n}u^2_ndx\to \int_{\mathbb{R}^3}\phi^s_{u}u^2dx,\quad
\int_{\mathbb{R}^3}\phi^s_{u_n}u_n\varphi dx\to \int_{\mathbb{R}^3}\phi^s_{u}u\varphi dx, \forall \varphi\in H^s_r(\mathbb{R}^3).
\]
\end{lemma}

\begin{lemma}\label{lem3.5}
Let $u\in H^s_r(\mathbb{R}^3)$ be a weak solution of  \eqref{lab1.5},
then we can construct the Pohoz\v{a}ev manifold
$$
\mathcal{P}_c= \{u\in S_r(c) : P_\mu(u)=0 \},
$$
where
\begin{align*}
P_\mu(u)
&=s\Big(a+b\int_{\mathbb{R}^3}|(-\Delta)^{s/2}u|^2dx\Big)
 \int_{\mathbb{R}^3}|(-\Delta)^{s/2}u|^2dx
  +\frac{3-2s}{4}\int_{\mathbb{R}^3}\phi^s_u u^2dx\\
&-s\mu \delta_{p,s}\int_{\mathbb{R}^3}|u|^pdx-s\int_{\mathbb{R}^3}|u|^{2^*_s}dx.
\end{align*}
\end{lemma}

\begin{proof}
From Proposition \ref{pro2.2}, we know that $u$ satisfies the Pohoz\v{a}ev identity
\begin{equation}\label{lab3.51}
\begin{split}
&\frac{3-2s}{2}\Big(a+b\int_{\mathbb{R}^3}|(-\Delta)^{s/2}u|^2dx\Big)
 \int_{\mathbb{R}^3}|(-\Delta)^{s/2}u|^2dx
 +\frac{3+2s}{4}\int_{\mathbb{R}^3}\phi^s_u u^2dx\\
&=\frac{3}{2}\lambda \int_{\mathbb{R}^3}|u|^2dx
 + \frac{3}{p}\mu \int_{\mathbb{R}^3}|u|^pdx
 +\frac{3}{2^*_s}\int_{\mathbb{R}^3}|u|^{2^*_s}dx.
\end{split}
\end{equation}
Moreover, since $u$ is the weak solution of equation ({\ref{lab1.5}}), we have
\begin{equation}\label{lab3.52}
\begin{split}
&\Big(a+b\int_{\mathbb{R}^3}|(-\Delta)^{s/2}u|^2dx\Big)
 \int_{\mathbb{R}^3}|(-\Delta)^{s/2}u|^2dx
 +\int_{\mathbb{R}^3}\phi^s_u u^2dx\\
&=\lambda \int_{\mathbb{R}^3}|u|^2dx+\mu \int_{\mathbb{R}^3}|u|^pdx
 +\int_{\mathbb{R}^3}|u|^{2^*_s}dx.
\end{split}
\end{equation}
Combining this with \eqref{lab3.51} and \eqref{lab3.52}, we obtain
\begin{align*}
&s\Big(a+b\int_{\mathbb{R}^3}|(-\Delta)^{s/2}u|^2dx\Big)
 \int_{\mathbb{R}^3}|(-\Delta)^{s/2}u|^2dx
 +\frac{3-2s}{4}\int_{\mathbb{R}^3}\phi^s_u u^2dx\\
&-s\mu \delta_{p,s}\int_{\mathbb{R}^3}|u|^pdx-s\int_{\mathbb{R}^3}|u|^{2^*_s}dx=0,
\end{align*}
which completes the proof.
\end{proof}

\section{Proof of Theorem \ref{thm1.1}}
In this section, we show the multiplicity of normalized solutions to  system \eqref{lab1.1}.
To begin by recalling the definition of genus.
Let $X$ be a Banach space and let $A$ be a subset of $X$. The set $A$ is said to be
symmetric if $u\in A$ implies that $-u\in A$. We denote the set
$$
\Sigma :=\{A\subset X\setminus \{0\}: A \text{ is closed and symmetric with respect
 to the origin}\}.
$$
For $A\in\Sigma$, define
\[
\gamma(A) =
\begin{cases}
0, & \text{if } A = \emptyset, \\
\inf\{k \in \mathbb{N}: \exists \text{ an odd }
\varphi \in C(A, \mathbb{R}^k \setminus \{0\})\}, \\
+\infty,& \text{if no such odd map exists,}
\end{cases}
\]
and that $\Sigma_k=\{A\in \Sigma:\gamma(A)\geq k\}$.

\begin{lemma}[\cite{Z1688}] \label{lem2.131}
Let $\{u_n\}$ be a bounded sequence in $D^{s,2}(\mathbb{R}^3)$ converging weakly
and a.e.\ to some $u\in D^{s,2}(\mathbb{R}^3)$.
 We have that $|(-\Delta)^{s/2}u_n|^2\rightharpoonup \omega$ and
 $|u_n|^{2^*_s}\rightharpoonup \zeta$ in the sense of measures.
Then, there exist some at most a countable set $J$, a family of points
$\{x_j\}_{j\in J} \subset \mathbb{R}^3$, and families of positive numbers
$\{\zeta_j\}_{j\in J}$ and $\{\omega_j\}_{j\in J}$ such that
\begin{gather}\label{lab3.11}
\omega \geq |(-\Delta)^{s/2}u|^2+\sum_{j\in J}\omega_j\delta_{x_j}, \\
\label{lab3.12}
\zeta=|u|^{2^*_s}+\sum_{j\in J}\zeta_j\delta_{x_j}, \\
\label{lab3.13}
\omega(\{x_j\})\geq S\zeta^{2/2^*_s}_j,
\end{gather}
where $\delta_{x_j}$ is the Dirac-mass of mass 1 concentrated at
$x_j\in {\mathbb{R}^3}$.
\end{lemma}


\begin{lemma}[\cite{Z1688}] \label{lem2.11}
Let $\{u_n\}\subset D^{s,2}(\mathbb{R}^3)$ be a sequence in Lemama \ref{lem2.131}
and define
\[
\omega_\infty:={\lim_{R\to \infty}}\limsup_{n\to\infty}
\int_{|x|\geq R} |(-\Delta)^{s/2}u_n|^2dx , \quad
\zeta_\infty:={\lim_{R\to \infty}}\limsup_{n\to\infty}
\int_{|x|\geq R} |u_n|^{2^*_s}dx.
\]
Then it follows that
\begin{gather*}%\label{lab2.111}
\omega_\infty \geq S\zeta^{2/2^*_s}_\infty, \\
%\label{lab2.112}
\limsup_{n\to\infty} \int_{\mathbb{R}^3}|(-\Delta)^{s/2}u_n|^2dx
=\int_{\mathbb{R}^3}d\omega +\omega_\infty,\\
%\label{lab2.113}
\limsup_{n\to\infty} \int_{\mathbb{R}^3}|u_n|^{2^*_s}dx=\int_{\mathbb{R}^3}d\zeta +\zeta_\infty.\\
\end{gather*}
\end{lemma}

For $u\in S_{r}(c)$, in view of Lemma \ref{lem3.2}, and the Sobolev inequality, one has
\begin{align*}
I_\mu(u)
&= \frac{a}{2}\int_{\mathbb{R}^3}|(-\Delta)^{s/2}u|^2dx
 +\frac{b}{4}\Big(\int_{\mathbb{R}^3}|(-\Delta)^
{\frac{s}{2}}u|^2dx\Big)^2
 +\frac{1}{4}\int_{\mathbb{R}^3}\phi^s_u u^2dx\\
&\quad -\frac{\mu}{p}\int_{\mathbb{R}^3}|u|^pdx
 -\frac{1}{2^*_s}\int_{\mathbb{R}^3}|u|^{2^*_s}dx\\
&\geq  \frac{a}{2}\|(-\Delta)^{s/2}u\|^2_2+\frac{b}{4}\|(-\Delta)^{s/2}u\|^4_2
 -\frac{\mu}{p}c^{p(1-\delta_{p,s})}C({p,s})\|(-\Delta)^{s/2}u\|^{p\delta_{p,s}}_2\\
&\quad -\frac{1}{2^*_s}S^{-\frac{2^*_s}{2}}\|(-\Delta)^{s/2}u\|^{2^*_s}_2\\
&:= h(\|(-\Delta)^{s/2}u\|_2),
\end{align*}
where
$$
h(r)=\frac{a}{2}r^2+\frac{b}{4}r^4
-\frac{\mu}{p}c^{p(1-\delta_{p,s})}C({p,s})r^{p\delta_{p,s}}
-\frac{1}{2^*_s}S^{-\frac{2^*_s}{2}}r^{2^*_s}.
$$
Recalling that $p\in(2,\frac{4s}{3}+2)$, we obtain that $p\delta_{p,s}<2$, and there
exists $\beta>0$ such that, if $\mu c^{p(1-\delta_{p,s})}<\beta$, the function
$h(\cdot)$ attains its positive local maximum. More precisely, there exist two
constants $0<R_1<R_2<+\infty$, such that
$$
h(r)>0, \forall r\in(R_1, R_2),\quad
 h(r)<0, \quad \forall r\in(0,R_1)\cup(R_2, +\infty).
$$
Let $\tau: \mathbb{R}^+\to [0,1]$ be a non-increasing and $C^\infty$ function
satisfying
\[
\tau(r) =
\begin{cases}
1, & \text{if } r\in [0, R_1], \\
0, & \text{if } r\in [R_2, +\infty).
\end{cases}
\]

In the sequel, we consider the truncated functional
\begin{equation}
\begin{split}
I_{\mu,\tau}(u)
&= \frac{a}{2}\int_{\mathbb{R}^3}|(-\Delta)^{s/2}u|^2dx+\frac{b}{4}
 \Big(\int_{\mathbb{R}^3}|(-\Delta)^{s/2}u|^2dx\Big)^2
 +\frac{1}{4}\int_{\mathbb{R}^3}\phi^s_u u^2dx\\
&\quad -\frac{\mu}{p}\int_{\mathbb{R}^3}|u|^pdx
 -\frac{\tau(\|(-\Delta)^{s/2}u\|_2)}{2^*_s}\int_{\mathbb{R}^3}|u|^{2^*_s}dx.
\end{split}
\end{equation}
For $u\in S_{r}(c)$, again by Lemma \ref{lem3.2}, and the Sobolev inequality,
it is easy to see that
\begin{align*}
I_{\mu,\tau}(u)
&\geq  \frac{a}{2}\|(-\Delta)^{s/2}u\|^2_2+\frac{b}{4}\|(-\Delta)^{s/2}u\|^4_2
 -\frac{\mu}{p}c^{p(1-\delta_{p,s})}C({p,s})\|(-\Delta)^{s/2}u\|^{p\delta_{p,s}}_2\\
&\quad -\frac{\tau(\|(-\Delta)^{s/2}u\|_2)}{2^*_s}\|(-\Delta)^{s/2}u\|^{2^*_s}_2\\
&:= \tilde{h}(\|(-\Delta)^{s/2}u\|_2),
\end{align*}
where
$$
\tilde{h}(r)=\frac{a}{2}r^2+\frac{b}{4}r^4
-\frac{\mu}{p}c^{p(1-\delta_{p,s})}C({p,s})r^{p\delta_{p,s}}
-\frac{\tau(r)}{2^*_s}S^{-\frac{2^*_s}{2}}r^{2^*_s}.
$$
Then, by the definition of $\tau(\cdot)$, when
$c\in \big(0, (\frac{\beta}{\mu})^{\frac{1}{p(1-\delta_{p, s})}}\big)$, we have
$$
\tilde{h}(r)<0, \forall r\in(0, R_1),\quad
\tilde{h}(r)>0, \forall r\in(R_1, +\infty).
$$
In what follows, we  assume that
$c\in \big(0, (\frac{\beta}{\mu})^{\frac{1}{p(1-\delta_{p, s})}}\big)$.
Without loss of generality, in the sequel, we  assume that
\begin{equation}\label{1.9}
\frac{a}{2}r^2+\frac{b}{4}r^4-\frac{1}{2^*_s}S^{-\frac{2^*_s}{2}}r^{2^*_s}\geq 0,
\quad \forall ~r\in [0, R_1]
\end{equation}

\begin{lemma}\label{lem2.161}
The functional $I_{\mu,\tau}$ has the following characteristics:
\begin{itemize}
\item[(i)] $I_{\mu,\tau} \in C^1(H^s_{r}(\mathbb{R}^3), \mathbb{R})$;

\item[(ii)] $I_{\mu,\tau}$ is coercive and bounded from below on $S_{r}(c)$.
Moreover, if $I_{\mu,\tau}(u)\leq 0$, then $\|(-\Delta)^{s/2}u\|_2\leq R_1$ and
$I_{\mu,\tau}(u)=I_{\mu}(u)$;

\item[(iii)] $ I_{\mu,\tau}|_{S_{r}(c)}$ satisfies the $(PS)_d$ condition for all
$d<min\{0, \Lambda-\mu c^{p(1-\delta_{p,s})}D\}$, provided that $\mu>\mu^*_1>0$
large, where $\Lambda=(\frac{1}{2}-\frac{1}{2^*_s}) a^{\frac{3}{2s}}S^{\frac{3}{2s}}$,
$D=(\frac{1}{p}-\frac{1}{2^*_s})C(p,s)R_1^{p\delta_{p,s}}$.
\end{itemize}
\end{lemma}

\begin{proof}
The proofs of (i) and (ii) are easy. To prove item (iii), Let
$\{u_n\}$ be a $(PS)_d$ sequence of $I_{\mu, \tau}$ restricted to $S_{r}(c)$
with $d<min\{0, \Lambda-\mu c^{q(1-\delta_{p,s})}D\}$. By $(ii)$, we see that
$\|(-\Delta)^\frac{s}{2}u_n\|_2\leq R_1$ for large $n$, and thus $\{u_n\}$ is a
$(PS)_d$ sequence of $I_{\mu}|_{S_{r}(c)}$ with $d<min\{0,
\Lambda-\mu c^{p(1-\delta_{p,s})}D\}$; i.e.,
$I_{\mu}(u_n)\to d$ and $\|I_{\mu}|'_{S_{r}(c)}(u_n)\|\to 0$ as $n\to \infty$.
Then, $\{u_n\}$ is bounded in $H^s_{r}(\mathbb{R}^3)$. Therefore, up to a
subsequence, there exists $u\in H^s_r(\mathbb{R}^3)$ such that
$u_n\rightharpoonup u$ in $H^s_{r}(\mathbb{R}^3)$ and
$u_n\to u$ in $L^p(\mathbb{R}^3)$ for $2<p<2^*_s$ and
$u_n(x)\to u(x)$ a.e. on $\mathbb{R}^3$.
From $2<p<\frac{4s}{3}+2<2^*_s$ and Lemma \ref{lem2.171}, we infer to
$$
\lim_{n\to \infty}\int_{\mathbb{R}^3}|u_n|^pdx
=\int_{\mathbb{R}^3}|u|^pdx, \quad
\int_{\mathbb{R}^3}\phi^s_{u_n}{u_n}^2dx\to \int_{\mathbb{R}^3}\phi^s_u u^2dx.
$$
Moreover, we have that $u\not\equiv 0$. Indeed, assume by contradiction that,
$u\equiv0$, then $\lim_{n\to \infty}\int_{\mathbb{R}^3}|u_n|^pdx=0$.
From \eqref{1.9} and the definition of $I_{\mu, \tau}$, we infer that
\begin{align*}
0>d
&=\lim_{n\to \infty}I_{\mu, \tau}(u_n)=\lim_{n\to \infty}I_{\mu}(u_n)\\
&=\lim_{n\to \infty}\Big[\frac{a}{2}\|(-\Delta)^{s/2}u_n\|^2_2
 +\frac{b}{4}\|(-\Delta)^{s/2}u_n\|^4_2
 +\frac{1}{4}\int_{\mathbb{R}^3}\phi^s_{u_n}{u_n}^2dx\\
&\quad -\frac{\mu}{p}\int_{\mathbb{R}^3}|u_n|^pdx
 -\frac{1}{2^*_s}\int_{\mathbb{R}^3}|u_n|^{2^*_s}dx\Big]\\
&\geq \lim_{n\to \infty}\Big[\frac{a}{2}\|(-\Delta)^{s/2}u_n\|^2_2+\frac{b}{4}\|(-\Delta)^{s/2}u_n\|^4_2
\\
&\quad -\frac{\mu}{p}\int_{\mathbb{R}^3}|u_n|^pdx
 -\frac{1}{2^*_s}S^{-\frac{2^*_s}{2}}\|(-\Delta)^{s/2}u_n\|^{2^*_s}_2\Big]\\
&\geq -\frac{\mu}{p}\int_{\mathbb{R}^3}|u|^pdx=0,
\end{align*}
which is absurd. On the other hand, setting the function
 $\Theta(v): H^s_{r}(\mathbb{R}^3)\to \mathbb{R}$ by
$$
\Theta(v)=\frac{1}{2}\int_{\mathbb{R}^3}|v|^2dx,
$$
it follows that $S_r(c)=\Theta^{-1}(\{\frac{c^2}{2}\})$.
Then, by \cite[Proposition 5.12]{W6888}, there exists $\lambda_n\in {\mathbb{R}}$
such that
$$
\|I'_\mu(u_n)-\lambda_n\Theta'(u_n)\|\to 0,~~~ as~ n\to \infty.
$$
Hence, in $H^{-s}_{r}(\mathbb{R}^3)$, we have
\[  %\label{4.1}
\Big(a+b\int_{\mathbb{R}^3}|(-\Delta)^{s/2}u_n|^2dx\Big)
 (-\Delta)^su_n+\phi^s_{u_n}u_n
-\mu |u_n|^{p-2}u_n-|u_n|^{2^*_s-2}u_n=\lambda_n u_n+o(1),
\]
where $H^{-s}_{r}(\mathbb{R}^3)$ is the dual space of $H^{s}_{r}(\mathbb{R}^3)$.
Thus, we have for $\varphi \in H^s_{r}(\mathbb{R}^3)$, that
\begin{equation}\label{4.11}
\begin{split}
&\Big(a+b\int_{\mathbb{R}^3}|(-\Delta)^{s/2}u_n|^2dx\Big)
 \int_{\mathbb{R}^3}(-\Delta)^{s/2}u_n (-\Delta)^{s/2}\varphi dx
 +\int_{\mathbb{R}^3}\phi^s_{u_n}u_n\varphi dx \\
&\quad -\mu \int_{\mathbb{R}^3}|u_n|^{p-2}u_n\varphi dx
 -\int_{\mathbb{R}^3}|u_n|^{2^*_s-2}u_n\varphi dx\\
&=\lambda_n \int_{\mathbb{R}^3} u_n\varphi dx+o(1),
\end{split}
\end{equation}
and if we choose $\varphi=u_n$, we obtain
\begin{equation}\label{4.111}
\begin{split}
&\Big(a+b\int_{\mathbb{R}^3}|(-\Delta)^{s/2}u_n|^2dx\Big)
 \int_{\mathbb{R}^3}|(-\Delta)^{s/2}u_n|^2 dx
 +\int_{\mathbb{R}^3}\phi^s_{u_n}{u^2_n}dx
 -\mu \int_{\mathbb{R}^3}|u_n|^{p}dx
 -\int_{\mathbb{R}^3}|u_n|^{2^*_s} dx \\
&=\lambda_n \int_{\mathbb{R}^3} {u^2_n} dx+o(1) ,
\end{split}
\end{equation}
from \eqref{4.111}, and the boundedness of $\{u_n\}$ in $D^{s, 2}(\mathbb{R}^3)$,
we can deduce that $\{\lambda_n\}$ is bounded in $\mathbb{R}$.
Then we can assume that, up to subsequence, $\lambda_n\to \lambda$ for some
$\lambda \in {\mathbb{R}}$.

Next, we shall prove $u_n\to u$ in $L^{2^*_s}(\mathbb{R}^3)$ by using  the
concentration-compactness principle due to Lions {\cite{L7889}}.
Since $\|(-\Delta)^{s/2}u_n\|_2\leq {R}_1$, for $n$ large enough,
 by Lemma \ref{lem2.131}, there exist two positive measures, $\zeta$,
 $\omega\in \mathcal{M}(\mathbb{R}^3)$, such that
\begin{equation}\label{4.16}
|(-\Delta)^{s/2}u_n|^2\rightharpoonup \omega, |u_n|^{2^*_s}\rightharpoonup\zeta
\quad\text{in }\mathcal{M}(\mathbb{R}^3)
\end{equation}
as $n\to \infty$. Then, by Lemma \ref{lem2.131}, either $u_n\to u$ in
$L^{2^*_s}_{loc}(\mathbb{R}^3)$ or there exists a (at most countable) set of
distinct points $\{x_j\}_{j\in J}\subset \mathbb{R}^3$ and positive numbers
$\{\zeta_j\}_{j\in J}$ such that
$$
\zeta=|u|^{2^*_s}+\Sigma_{j\in J}\zeta_j \delta_{x_j}.
$$
Moreover, there exist some at most a countable set $J\subset \mathbb{N}$,
a corresponding set of distinct points $\{x_j\}_{j\in J}\subset \mathbb{R}^3$,
and two sets of positive numbers  $\{\zeta_j\}_{j\in J}\subset \mathbb{R}^3$ and
$\{\omega_j\}_{j\in J}\subset \mathbb{R}^3$ such that items
\eqref{lab3.11}-\eqref{lab3.13} hold. Now, assume that $J\neq \emptyset$.
We split the proof into three steps.
\smallskip

\noindent\textbf{Step 1.} We prove that $a\omega(\{x_j\})\leq\zeta_j$,
where $\omega(\{x_j\})$, and $\zeta_j$ come from Lemma \ref{lem2.131}.
We define $\varphi \in C^\infty_0(\mathbb{R}^3)$ as a cut-off function with
$\varphi \in [0,1]$, $\varphi=1$ in $B_{\frac{1}{2}}(0)$, $\varphi=0$ in
$\mathbb{R}^3\setminus B_1(0)$. For any $\rho>0$, define
\[
\varphi_\rho(x):=\varphi\Big(\frac{x-x_j}{\rho}\Big)
=\begin{cases}
1, & |x-x_j|\leq \frac{\rho}{2}, \\
0, & |x-x_j|\geq {\rho}.
\end{cases}
\]
By the boundedness of $\{u_n\}$ in $H^s_{r}(\mathbb{R}^3)$, we have that
$\{u_n\varphi_\rho\}$ is also bounded in $H^s_{r}(\mathbb{R}^3)$. Thus, one has
\begin{equation}\label{4.17}
\begin{split}
o_n(1)
&= \langle I'_\mu(u_n), u_n\varphi_\rho \rangle\\
&= \Big(a+b\int_{\mathbb{R}^3}|(-\Delta)^{s/2}u_n|^2dx\Big)
 \int_{\mathbb{R}^3}(-\Delta)^{s/2}u_n
 (-\Delta)^{s/2}(u_n\varphi_\rho)dx+\int_{\mathbb{R}^3}\phi^s_{u_n}
 u_n \varphi_\rho dx\\
&\quad-\mu \int_{\mathbb{R}^3}|u_n|^p\varphi_\rho dx
 -\int_{\mathbb{R}^3}|u_n|^{2^*_s}\varphi_\rho dx.
\end{split}
\end{equation}
It is easy to check that
\begin{align*}
\int_{\mathbb{R}^3}(-\Delta)^{s/2}u_n
(-\Delta)^{s/2}(u_n\varphi_\rho)dx 
&= \int\int_{\mathbb{R}^6}\frac{(u_n(x)-u_n(y))(u_n(x)\varphi_\rho(x)-u_n(y)\varphi_\rho(y))}{|x-y|^{3+2s}}\,dx\,dy\\
&= \int\int_{\mathbb{R}^6}\frac{|u_n(x)-u_n(y)|^2\varphi_\rho(y)}{|x-y|^{3+2s}}\,dx\,dy\\
&\quad+\int\int_{\mathbb{R}^6}\frac{(u_n(x)-u_n(y))(\varphi_\rho(x)
 -\varphi_\rho(y))u_n(x)}{|x-y|^{3+2s}}\,dx\,dy \\
&=T_1+T_2\,.
\end{align*}
For $T_1$, by \eqref{4.16}, we obtain
\[
\begin{split}
\lim_{\rho\to 0}\lim_{n\to\infty}T_1
&= \lim_{\rho\to 0}\lim_{n\to\infty}\int\int_{\mathbb{R}^6}\frac{|u_n(x)-u_n(y)|^2\varphi_\rho(y)}{|x-y|^{3+2s}}\,dx\,dy\\
&= \lim_{\rho\to 0}\int_{\mathbb{R}^3}\varphi_\rho d\omega=\omega(\{x_j\}).
\end{split}
\]
From H\"older's inequality, we have
\[
\begin{split}
|T_2|
&= \Big|\int\int_{\mathbb{R}^6}\frac{(u_n(x)-u_n(y))(\varphi_\rho(x)
 -\varphi_\rho(x))u_n(x)}{|x-y|^{3+2s}}\,dx\,dy\Big|\\
&\leq \int\int_{\mathbb{R}^6}
 \Big|\frac{(u_n(x)-u_n(y))(\varphi_\rho(x)-\varphi_\rho(x))u_n(x)}{|x-y|^{3+2s}}
 \Big|\,dx\,dy\\
&\leq \Big(\int\int_{\mathbb{R}^6}
 \frac{|\varphi_\rho(x)-\varphi_\rho(x)|^2u^2_{n}(x)}{|x-y|^{3+2s}}\,dx\,dy
  \Big)^{1/2}
\Big(\int\int_{\mathbb{R}^6}\frac{|u_n(x)-u_n(y)|^2}{|x-y|^{3+2s}}\,dx\,dy\Big)^{1/2}\\
&\leq C_1\Big(\int\int_{\mathbb{R}^6}\frac{|\varphi_\rho(x)-\varphi_\rho(x)
 |^2u^2_{n}(x)}{|x-y|^{3+2s}}\,dx\,dy\Big)^{1/2}.
\end{split}
\]
Analogously to the proof of \cite[Lemma 3.4]{Z1688}, one has
\begin{equation}
\begin{gathered}
\lim_{\rho\to 0}\lim_{n\to\infty}\int\int_{\mathbb{R}^6}\frac{|\varphi_\rho(x)-\varphi_\rho(y)|^2u^2_{n}(x)}
{|x-y|^{3+2s}}\,dx\,dy=0, \\
\label{4.18}
\lim_{\rho\to 0}\lim_{n\to\infty}\int_{\mathbb{R}^3}(-\Delta)^{s/2}u_n
(-\Delta)^{s/2}(u_n\varphi_\rho)dx=\omega(\{x_j\}).
\end{gathered}
\end{equation}
Again by \eqref{4.16}, we have
\begin{equation}\label{4.20}
\lim_{\rho\to 0}\lim_{n\to\infty}\int_{\mathbb{R}^3}|u_n|^{2^*_s}\varphi_\rho dx
=\lim_{\rho\to 0}\int_{\mathbb{R}^3}\varphi_\rho d\zeta=\zeta(\{x_j\})=\zeta_j.
\end{equation}
By the definition of $\varphi_\rho$, and the absolute continuity of the Lebesgue
integral, one has
\begin{equation}\label{4.21}
\lim_{\rho\to 0}\lim_{n\to\infty}\int_{\mathbb{R}^3}|u_n|^{p}\varphi_\rho dx
=\lim_{\rho\to 0}\int_{\mathbb{R}^3}|u|^{p}\varphi_\rho dx
=\lim_{\rho\to 0}\int_{|x-x_j|\leq\rho}|u|^{p}\varphi_\rho dx=0.
\end{equation}
Thus, by Proposition {\ref{pro2.1}} and Lemma \ref{lem2.1311}, we have
\begin{align*} %\label{4.22}
\int_{\mathbb{R}^3}\phi^s_{u_n}u^2_n\varphi_\rho dx&\leq
C\left(\int_{\mathbb{R}^3}|u_n|^{\frac{12}{3+2s}}dx\right)^{\frac{3+2s}{6}}\left(\int_{\mathbb{R}^3}|u^2_n\varphi_\rho|^{\frac{6}{3+2s}}dx\right)
^{\frac{3+2s}{6}}\\
&\leq  C_2\|u_n\|^2_{H^s_r}
 \Big(\int_{\mathbb{R}^3}{|u_n|}^{\frac{12}{3+2s}}|\varphi_\rho|^{\frac{6}{3+2s}}dx
  \Big)^{\frac{3+2s}{6}}\\
&\leq  C_3 \Big(\int_{\mathbb{R}^3}{|u_n|}^{\frac{12}{3+2s}}\varphi_\rho dx
 \Big)^{\frac{3+2s}{6}}.
\end{align*}
Therefore,
\begin{equation}\label{4.23}
\begin{split}
\lim_{\rho\to 0}\lim_{n\to\infty}\int_{\mathbb{R}^3}\phi^s_{u_n}u^2_n\varphi_\rho dx
&\leq  \lim_{\rho\to 0}\lim_{n\to\infty} C_3
 \Big(\int_{\mathbb{R}^3}{|u_n|}^{\frac{12}{3+2s}}\varphi_\rho dx\Big)^{\frac{3+2s}{6}}\\
&= \lim_{\rho\to 0}C_3
 \Big(\int_{\mathbb{R}^3}{|u|}^{\frac{12}{3+2s}}\varphi_\rho dx\Big)^{\frac{3+2s}{6}}\\
&= \lim_{\rho\to 0}C_3
 \Big(\int_{|x-x_j|\leq\rho}{|u|}^{\frac{12}{3+2s}}\varphi_\rho dx\Big)^{\frac{3+2s}{6}}=0.
\end{split}
\end{equation}
Summing \eqref{4.17}-\eqref{4.23}, taking the limit as $n\to\infty$, and
then the limit as $\rho\to 0$, we arrive at
$$
\zeta_j\geq a\omega(\{x_j\}).
$$
\smallskip

\noindent\textbf{Step 2.}
We show that $a\omega_\infty\leq\zeta_\infty$, where $\omega_\infty$ and
$\zeta_\infty$ are given in Lemma \ref{lem2.11}.
Let $\psi \in C^\infty_0(\mathbb{R}^3)$ be a cut-off function with $\psi \in[0,1]$,
$\psi=0$ in $B_{\frac{1}{2}}(0)$, $\psi=1$ in $\mathbb{R}^3\setminus B_1(0)$.
For any $R>0$, define
\[
\psi_R(x):=\psi\left(\frac{x}{R}\right) =
\begin{cases}
0, & |x|\leq \frac{R}{2}, \\
1, & |x|\geq {R}.
\end{cases}
\]
Using again the boundedness of $\{u_n\}$ and $\{u_n\psi_R\}$ in $H^s_{r}(\mathbb{R}^3)$,
we have
\begin{equation}\label{4.24}
\begin{split}
o_n(1)
&= \langle I'_\mu(u_n), u_n\psi_R \rangle\\
&= \Big(a+b\int_{\mathbb{R}^3}|(-\Delta)^{s/2}u_n|^2dx\Big)
 \int_{\mathbb{R}^3}(-\Delta)^{s/2}u_n(-\Delta)^{s/2}(u_n\psi_R)dx\\
&\quad +\int_{\mathbb{R}^3}\phi^s_{u_n} u^2_n\psi_R dx
 -\mu \int_{\mathbb{R}^3}|u_n|^p\psi_R  dx-\int_{\mathbb{R}^3}|u_n|^{2^*_s}\psi_R  dx.
\end{split}
\end{equation}
It is easy to derive that
\[
\begin{split}
&\int_{\mathbb{R}^3}(-\Delta)^{s/2}u_n
(-\Delta)^{s/2}(u_n\psi_R)dx\\
&= \int\int_{\mathbb{R}^6}\frac{(u_n(x)-u_n(y))(u_n(x)\psi_R(x)-u_n(y)\psi_R(y))}{|x-y|^{3+2s}}\,dx\,dy\\
&= \int\int_{\mathbb{R}^6}\frac{|u_n(x)-u_n(y)|^2\psi_R(y)}{|x-y|^{3+2s}}\,dx\,dy\\
&\quad+\int\int_{\mathbb{R}^6}\frac{(u_n(x)-u_n(y))(\psi_R(x)-\psi_R(y))u_n(x)}{|x-y|^{3+2s}}\,dx\,dy
=T_3+T_4\,.
\end{split}
\]
For $T_3$, by \eqref{4.16} and Lemma \ref{lem2.11}, we infer that
\[
\lim_{R\to \infty}\lim_{n\to\infty}T_3
= \lim_{R\to \infty}\lim_{n\to\infty}\int\int_{\mathbb{R}^6}
\frac{|u_n(x)-u_n(y)|^2\psi_R(y)}{|x-y|^{3+2s}}\,dx\,dy
=\omega_\infty.
\]
From H\"older's inequality, we have
\begin{align*}
|T_4|
&= \Big|\int\int_{\mathbb{R}^6}\frac{(u_n(x)-u_n(y))(\psi_R(x)
 -\psi_R(y))u_n(x)}{|x-y|^{3+2s}}\,dx\,dy\Big|\\
&\leq \int\int_{\mathbb{R}^6}
 \Big|\frac{(u_n(x)-u_n(y))(\psi_R(x)-\psi_R(y))u_n(x)}{|x-y|^{3+2s}}\Big|\,dx\,dy\\
&\leq \Big(\int\int_{\mathbb{R}^6}\frac{|\psi_R(x)-\psi_R(x)|^2|u_{n}(x)|^2}{|x-y|^{3+2s}}
 \,dx\,dy\Big)^{1/2}
\Big(\int\int_{\mathbb{R}^6}\frac{|u_n(x)-u_n(y)|^2}{|x-y|^{3+2s}}\,dx\,dy\Big)^{1/2}\\
&\leq C\Big(\int\int_{\mathbb{R}^6}\frac{|\psi_R(x)
 -\psi_R(y)|^2|u_{n}(x)|^2}{|x-y|^{3+2s}}\,dx\,dy\Big)^{1/2}.
\end{align*}
Combining the above, we conclude that
\begin{align*}
&\lim_{R\to \infty}\lim_{n\to\infty}\int\int_{\mathbb{R}^6}\frac{|\psi_R(x)
 -\psi_R(y)|^2|u_{n}(x)|^2}{|x-y|^{3+2s}}\,dx\,dy\\
&=\lim_{R\to \infty}\lim_{n\to\infty}\int\int_{\mathbb{R}^6}
 \frac{|[1-\psi_R(x)]-[1-\psi_R(y)]|^2|u_{n}(x)|^2}{|x-y|^{3+2s}}\,dx\,dy=0.
\end{align*}
Hence,
\begin{equation}\label{4.244}
\lim_{R\to \infty}\lim_{n\to\infty}\int_{\mathbb{R}^3}(-\Delta)^{s/2}u_n
(-\Delta)^{s/2}(u_n\psi_R)dx=\omega_\infty.
\end{equation}
By Lemma \ref{lem2.11}, one has
\begin{equation}\label{4.25}
\lim_{R\to \infty}\lim_{n\to\infty}\int_{\mathbb{R}^3}|u_n|^{2^*_s}\psi_R dx
=\zeta_\infty.
\end{equation}
Analogous the proof of \cite[Lemma 3.3]{Z1688}, we infer that
\begin{equation}\label{4.26}
\lim_{R\to \infty}\lim_{n\to\infty}\int_{\mathbb{R}^3}|u_n|^{p}\psi_R dx
=\lim_{R\to \infty}\int_{\mathbb{R}^3}|u|^{p}\psi_R dx
=\lim_{R\to \infty}\int_{|x|\geq\frac{R}{2}}|u|^{p}\psi_R dx=0.
\end{equation}
Moreover, we can obtain
\begin{equation}\label{4.27}
\begin{split}
\lim_{R\to \infty}\lim_{n\to\infty}\int_{\mathbb{R}^3}\phi^s_{u_n}u^2_n\psi_R dx
&\leq  \lim_{R\to \infty}\lim_{n\to\infty} C_3 
 \Big(\int_{\mathbb{R}^3}{|u_n|}^{\frac{12}{3+2s}}\psi_Rdx\Big)^{\frac{3+2s}{6}}\\
&= \lim_{R\to \infty}C_3 
 \Big(\int_{\mathbb{R}^3}{|u|}^{\frac{12}{3+2s}}\psi_Rdx\Big)^{\frac{3+2s}{6}}\\
&= \lim_{R\to \infty}C_3
 \Big(\int_{{|x|\geq\frac{R}{2}}}{|u|}^{\frac{12}{3+2s}}\psi_R dx\Big)^{\frac{3+2s}{6}}=0.
\end{split}
\end{equation}
Summing up, from \eqref{4.24}-\eqref{4.27}, taking the limit as $n\to \infty$,
and then the limit as $R\to\infty$, we have
$$
\zeta_\infty \geq a\omega_\infty.
$$
\smallskip

\noindent\textbf{Step 3.}
 We claim that $\zeta_j=0$ for any $j\in J$ and $\zeta_\infty=0$.
Suppose by contradiction that, there exists $j_0\in J$ such that $\zeta_{j_0}>0$
or $\zeta_\infty>0$. Then step 1, step 2,  Lemma \ref{lem2.131} and
Lemma \ref{lem2.11} imply that
\[  %\label{4.28}
\zeta_{j_0}\leq (S^{-1}{\omega{(\{x_{j_0}\}))}}^{\frac{2^*_s}{2}}
\leq(S^{-1}a^{-1}\zeta_{j_0})^{\frac{2^*_s}{2}},\quad
\zeta_\infty \leq (S^{-1}{\omega_\infty)}^{\frac{2^*_s}{2}}
\leq(S^{-1}a^{-1}\zeta_\infty)^{\frac{2^*_s}{2}}.
\]
Consequently, we obtain $\zeta_{j_0}\geq (aS)^{\frac{3}{2s}}$ or
$\zeta_{\infty}\geq (aS)^{\frac{3}{2s}}$.
If $\zeta_{j_0}\geq (aS)^{\frac{3}{2s}}$, one has
\begin{align*}
d&= \lim_{n\to \infty}\Big[I_\mu(u_n)-\frac{1}{2^*_s}
 \langle I'_\mu(u_n), u_n\rangle\Big]\\
&=  \lim_{n\to \infty}\Big[\Big(\frac{1}{2}-\frac{1}{2^*_s}\Big)
 a\|(-\Delta)^{s/2}u_n\|^2_2+\Big(\frac{1}{4}-\frac{1}{2^*_s}\Big)
  b\|(-\Delta)^{s/2}u_n\|^4_2
\\
&\quad +\Big(\frac{1}{4}-\frac{1}{2^*_s}\Big)
 \int_{\mathbb{R}^3}\phi^s_{u_n}u^2_ndx
 -\Big(\frac{1}{p}-\frac{1}{2^*_s}\Big)\mu \int_{\mathbb{R}^3}|u_n|^pdx\Big]\\
&\geq  \lim_{n\to \infty}\Big[\Big(\frac{1}{2}-\frac{1}{2^*_s}\Big)a
 \|(-\Delta)^{s/2}u_n\|^2_2
 -\Big(\frac{1}{p}-\frac{1}{2^*_s}\Big)\mu \int_{\mathbb{R}^3}|u_n|^pdx\Big]\\
&\geq \lim_{n\to \infty}\Big(\frac{1}{2}-\frac{1}{2^*_s}\Big)aS\|u_n\|^2_{2^*_s}
 -\lim_{n\to \infty}\Big(\frac{1}{p}-\frac{1}{2^*_s}\Big)
  \mu c^{p(1-\delta_{p,s})}C(p,s)\|(-\Delta)^{s/2}u_n\|^{p\delta_{p,s}}_2\\
&\geq \Big(\frac{1}{2}-\frac{1}{2^*_s}\Big)
 aS\lim_{\rho\to 0}\lim_{n\to \infty}
 \Big(\int_{\mathbb{R}^3}|u_n|^{2^*_s}\varphi_\rho dx\Big)^{2/2^*_s}
  -\Big(\frac{1}{p}-\frac{1}{2^*_s}\Big)\mu c^{p(1
  -\delta_{p,s})}C(p,s)R_1^{p\delta_{p,s}}\\
&=  \Big(\frac{1}{2}-\frac{1}{2^*_s}\Big)aS(\zeta_{j_0})^\frac{2}{2^*_s}
 -\Big(\frac{1}{p}-\frac{1}{2^*_s}\Big)\mu c^{p(1-\delta_{p,s})}C(p,s)R_1^{p\delta_{p,s}}\\
&\geq  \Big(\frac{1}{2}-\frac{1}{2^*_s}\Big) a^{\frac{3}{2s}}S^{\frac{3}{2s}}
 -\big(\frac{1}{p}-\frac{1}{2^*_s}\Big)\mu c^{p(1-\delta_{p,s})}C(p,s)R_1^{p\delta_{p,s}}\\
&= \Lambda-\mu c^{p(1-\delta_{p,s})}D,
\end{align*}
which contradicts $d<\min\{0, \Lambda-\mu c^{p(1-\delta_{p,s})}D\}$.
If $\zeta_{\infty}\geq (aS)^{\frac{3}{2s}}$, we have
\begin{align*}
d&= \lim_{n\to \infty}\Big[I_\mu(u_n)-\frac{1}{2^*_s}\langle I'_\mu(u_n),
 u_n\rangle\Big]\\
&=  \lim_{n\to \infty}\Big[\Big(\frac{1}{2}-\frac{1}{2^*_s}\Big)a\|
(-\Delta)^{s/2}u_n\|^2_2+\Big(\frac{1}{4}-\frac{1}{2^*_s}\Big)b\|(-\Delta)^{s/2}u_n\|^4_2
\\
&+\Big(\frac{1}{4}-\frac{1}{2^*_s}\Big)\int_{\mathbb{R}^3}\phi^s_{u_n}u^2_ndx
 -\Big(\frac{1}{p}-\frac{1}{2^*_s}\Big)\mu \int_{\mathbb{R}^3}|u_n|^pdx\Big]\\
&\geq  \lim_{n\to \infty}\Big[\Big(\frac{1}{2}-\frac{1}{2^*_s}\Big)a
\|(-\Delta)^{s/2}u_n\|^2_2-\Big(\frac{1}{p}-\frac{1}{2^*_s}\Big)
 \mu \int_{\mathbb{R}^3}|u_n|^pdx\Big]\\
&\geq \lim_{n\to \infty}\Big(\frac{1}{2}-\frac{1}{2^*_s}\Big)aS\|u_n\|^2_{2^*_s}
 -{\lim_{n\to \infty}}\Big(\frac{1}{p}-\frac{1}{2^*_s}\Big)
 \mu c^{p(1-\delta_{p,s})}C(p,s)\|(-\Delta)^{s/2}u_n\|^{p\delta_{p,s}}_2\\
&\geq \Big(\frac{1}{2}-\frac{1}{2^*_s}\Big)aS\lim_{R\to \infty}
 \lim_{n\to \infty}\Big(\int_{\mathbb{R}^3}|u_n|^{2^*_s}\psi_R dx\Big)^{2/2^*_s}
  -\Big(\frac{1}{p}-\frac{1}{2^*_s}\Big)\mu c^{p(1-\delta_{p,s})}C(p,s)
  R_1^{p\delta_{p,s}}\\
&= \Big(\frac{1}{2}-\frac{1}{2^*_s}\Big)aS(\zeta_\infty)^{2/2^*_s}
 -\Big(\frac{1}{p}-\frac{1}{2^*_s}\Big)\mu c^{p(1-\delta_{p,s})}C(p,s)
  R_1^{p\delta_{p,s}}\\
&\geq  \Big(\frac{1}{2}-\frac{1}{2^*_s}\Big) a^{\frac{3}{2s}}S^{\frac{3}{2s}}
 -\Big(\frac{1}{p}-\frac{1}{2^*_s}\Big)\mu c^{p(1-\delta_{p,s})}C(p,s)
 R_1^{p\delta_{p,s}}\\
&= \Lambda-\mu c^{p(1-\delta_{p,s})}D,
\end{align*}
which also contradicts $d<min\{0, \Lambda-\mu c^{p(1-\delta_{p,s})}D\}$.
Therefore, $\zeta_j=0$ for any $j\in J$ and $\zeta_\infty=0$. As a result,
by Lemma \ref{lem2.131}, we obtain that $u_n\to u$ in $L^{2^*_s}_{loc}(\mathbb{R}^3)$.
Combining with Lemma \ref{lem2.11}, we know that $u_n\to u$ in
$L^{2^*_s}(\mathbb{R}^3)$.

Now, we prove there exists $\mu^*_1>0$ independently on $n\in \mathbb{N}$ such that if 
$\mu>\mu^*_1$, the Lagrange multiplier $\lambda<0$. Indeed, note that 
$\{u_n\}\subset S_{r}(c)$ and $\|(-\Delta)^{s/2}u_n\|_2\leq R_1$ for large $n$, 
as can be seen from the previous proof of this Lemma, and \eqref{3.3}-\eqref{3.2} that, 
there exists $Q_1>0$ independently on $n$, such that
for large $n$
\begin{gather}\label{4.32}
\begin{split}
Q_1
&\leq \int_{\mathbb{R}^3}|u_n|^pdx \\
&\leq  C(p,s) \|(-\Delta)^{s/2}u_n\|^{p\delta_{p,s}}_2\|u_n\|^{p(1-\delta_{p,s})}_2\\
&\leq  C(p,s) {R_1}^{p\delta_{p,s}}c^{p(1-\delta_{p,s})},
\end{split} \\
\label{4.33}
\begin{split}
\int_{\mathbb{R}^3}\phi^s_{u_n} u^2_ndx\leq \Gamma_s \|u_n\|^4_{\frac{12}{3+2s}}
&\leq  \Gamma_sC(p_s,s)^{\frac{3+2s}{3}}
\|(-\Delta)^{s/2}u_n\|^{\frac{3-2s}{s}}_2\|u_n\|^{\frac{6s-3}{s}}_2\\
&\leq \Gamma_sC(p_s,s)^{\frac{3+2s}{3}}
R^{\frac{3-2s}{s}}_1c^{\frac{6s-3}{s}}\\
&:= Q_2,
\end{split}
\end{gather}
where $p_s:=\frac{12}{3+2s}$ and $Q_2=Q_2(s, R_1, c)>0 $. 
We define the constant
\begin{equation}\label{4.34}
\mu^*_1 =\frac{p(6s-3)Q_2}{2[6-p(3-2s)]Q_1}.
\end{equation}
By \eqref{4.32}-\eqref{4.34}, we have
\begin{equation}\label{4.35}
\mu^*_1 >\lim_{n\to +\infty}\frac{p(6s-3)\int_{\mathbb{R}^3}\phi^s_{u_n} u^2_ndx}
{2[6-p(3-2s)]\int_{\mathbb{R}^3}|u_n|^pdx}
=\frac{p(6s-3)\int_{\mathbb{R}^3}\phi^s_{u} u^2dx}{2[6-p(3-2s)]
\int_{\mathbb{R}^3}|u|^pdx}>0.
\end{equation}
Set $B:=\lim_{n\to \infty} \|(-\Delta)^{s/2}u_n\|^2_2\geq 0$, then, 
$0<\|(-\Delta)^{s/2}u\|^2_2\leq B$. For  any $\varphi \in H^s_{r}(\mathbb{R}^3)$, 
it follows by $u_n\rightharpoonup u$ in $H^s_{r}(\mathbb{R}^3)$ and 
$\lambda_n\to \lambda$, that
$$
\int_{\mathbb{R}^3}(-\Delta)^{s/2}u_n(-\Delta)^{s/2}\varphi dx
 \to \int_{\mathbb{R}^3}(-\Delta)^{s/2}u(-\Delta)^{s/2}\varphi dx 
 \quad \text{and}\quad
\lambda_n\int_{\mathbb{R}^3}u_n\varphi dx\to \lambda\int_{\mathbb{R}^3}u\varphi dx
$$
as $n\to \infty$. Since $\{|u_n|^{2^*_s-2}u_n\}$ is bounded in 
$L^{\frac{2^*_s}{2^*_s-1}}(\mathbb{R}^3)$,  $\{|u_n|^{p-2}u_n\}$ is bounded in 
$L^{\frac{2^*_s}{p-1}}(\mathbb{R}^3)$, and $u_n(x)\to u(x)$ a.e.\ in $\mathbb{R}^3$, 
we obtain that
$$
|u_n|^{2^*_s-2}u_n\rightharpoonup|u|^{2^*_s-2}u \quad \text{in } 
L^{\frac{2^*_s}{2^*_s-1}}(\mathbb{R}^3)\quad \text{and} \quad
|u_n|^{p-2}u_n\rightharpoonup|u|^{p-2}u \quad \text{in }L^{\frac{2^*_s}{p-1}}(\mathbb{R}^3),
$$
and so
$$
\int_{\mathbb{R}^3}|u_n|^{2^*_s-2}u_n\varphi dx
\to\int_{\mathbb{R}^3}|u|^{2^*_s-2}u\varphi dx \quad\text{and}\quad
\int_{\mathbb{R}^3}|u_n|^{p-2}u_n\varphi dx
\to\int_{\mathbb{R}^3}|u|^{p-2}u\varphi dx,
$$
as $n\to \infty$. Recall from Lemma \ref{lem2.171} that
$$
\int_{\mathbb{R}^3} \phi^s_{u_n}u_n\varphi dx
\to \int_{\mathbb{R}^3} \phi^s_{u}u\varphi dx, \quad
 \forall \varphi \in H^s_{r}(\mathbb{R}^3).
$$
Thus,  by \eqref{4.11}, for all $\varphi \in H^s_{r}(\mathbb{R}^3)$, we have
\begin{equation}\label{4.36}
\begin{split}
&(a+bB)\int_{\mathbb{R}^3}(-\Delta)^{s/2}u(-\Delta)^{s/2}\varphi dx+\int_{\mathbb{R}^3}\phi^s_{u}u\varphi dx\\
&-\mu \int_{\mathbb{R}^3}|u|^{p-2}u\varphi dx
 -\int_{\mathbb{R}^3}|u|^{2^*_s-2}u\varphi dx \\
&=\lambda \int_{\mathbb{R}^3} u\varphi dx,
\end{split}
\end{equation}
we can derive that $u$ solves the   equation
\begin{equation}\label{4.37}
(a+bB)(-\Delta)^su+\phi^s_{u}u-\mu |u|^{q-2}u-|u|^{2^*_s-2}u=\lambda u.
\end{equation}
 Moreover, by Lemma \ref{lem3.5}, $u$ satisfies
\begin{equation}\label{4.38}
s(a+bB)\int_{\mathbb{R}^3}|(-\Delta)^{s/2}u|^2dx+\frac{3-2s}{4}\int_{\mathbb{R}^3}
 \phi^s_u u^2dx
-s\mu \delta_{p,s}\int_{\mathbb{R}^3}|u|^pdx-s\int_{\mathbb{R}^3}|u|^{2^*_s}dx=0.
\end{equation}
Combining \eqref{4.37} and \eqref{4.38}, one has
\begin{equation}\label{4.39}
s\lambda\|u\|^2_2=\frac{6s-3}{4}\int_{\mathbb{R}^3}\phi^s_u u^2dx+s\mu \frac{(3-2s)p-6}{2p}\int_{\mathbb{R}^3}|u|^pdx.
\end{equation}
Now, if $\mu>\mu^*_1$, we conclude from \eqref{4.35}, that
\[
\mu >\frac{p(6s-3)\int_{\mathbb{R}^3}\phi^s_{u} u^2dx}{2[6-p(3-2s)]
\int_{\mathbb{R}^3}|u|^pdx}.
\]
Thus, from \eqref{4.39}, we infer to $\lim_{n\to+\infty}\lambda_n=\lambda<0$.
 By \eqref{4.111} and \eqref{4.36}, we derive
\begin{equation}\label{4.40}
\begin{split}
&\lim_{n\to\infty}\Big[\Big(a+b\int_{\mathbb{R}^3}|(-\Delta)^{s/2}u_n|^2dx\Big)
\int_{\mathbb{R}^3}|(-\Delta)^{s/2}u_n|^2 dx+\int_{\mathbb{R}^3}\phi^s_{u_n}{u^2_n}dx
-\lambda_n \int_{\mathbb{R}^3} {u^2_n} dx\Big]\\
&=\lim_{n\to\infty} \Big[\mu \int_{\mathbb{R}^3}|u_n|^{p}dx
 +\int_{\mathbb{R}^3}|u_n|^{2^*_s} dx\Big]\\
&=\mu \int_{\mathbb{R}^3}|u|^{p}dx+\int_{\mathbb{R}^3}|u|^{2^*_s} dx\\
&=(a+bB)\int_{\mathbb{R}^3}|(-\Delta)^{s/2}u|^2 dx+\int_{\mathbb{R}^3}\phi^s_{u}{u^2}dx
 -\lambda \int_{\mathbb{R}^3} {u^2} dx.
\end{split}
\end{equation}
Since $\lambda<0$ for $\mu>\mu^*_1$ large,  by Fatou's Lemma we obtain,
\begin{equation}\label{4.41}
\begin{split}
&\lim_{n\to\infty}\Big[\Big(a+b\int_{\mathbb{R}^3}|(-\Delta)^{s/2}u_n|^2dx\Big)
 \int_{\mathbb{R}^3}|(-\Delta)^{s/2}u_n|^2 dx
 +\int_{\mathbb{R}^3}\phi^s_{u_n}{u^2_n}dx-\lambda_n \int_{\mathbb{R}^3} {u^2_n} dx\Big]\\
&=\lim_{n\to\infty}\Big[(a+bB)\int_{\mathbb{R}^3}|(-\Delta)^{s/2}u_n|^2 dx
 +\int_{\mathbb{R}^3}\phi^s_{u_n}{u^2_n}dx-\lambda \int_{\mathbb{R}^3} {u^2_n} dx\Big]\\
&\geq (a+bB)\int_{\mathbb{R}^3}|(-\Delta)^{s/2}u|^2 dx
 +\int_{\mathbb{R}^3}\phi^s_{u}{u^2}dx-\liminf_{n\to\infty}\lambda 
 \int_{\mathbb{R}^3} {u^2_n} dx,
\end{split}
\end{equation}
and from \eqref{4.40}-\eqref{4.41}, one has
\begin{equation}\label{4.42}
-\lambda \int_{\mathbb{R}^3} {u^2} dx\geq\liminf_{n\to\infty}
 \Big(-\lambda \int_{\mathbb{R}^3} {u^2_n} dx\Big).
\end{equation}
But by Fatou's Lemma, we have
\begin{equation}\label{4.4300}
 \liminf_{n\to\infty}\Big(-\lambda\int_{\mathbb{R}^3} {u^2_n} dx\Big)
 \geq-\lambda \int_{\mathbb{R}^3} {u^2} dx.
\end{equation}
Combining \eqref{4.42} with \eqref{4.4300}, we obtain
\[ \label{4.43}
\lim_{n\to\infty}\Big(-\lambda \int_{\mathbb{R}^3} {u^2_n} dx\Big)
=-\lambda \int_{\mathbb{R}^3} {u^2} dx;
\]
that is,
\[ \label{4.43}
\lim_{n\to\infty} \int_{\mathbb{R}^3} {u^2_n} dx= \int_{\mathbb{R}^3} {u^2} dx.
\]
Thus, by \eqref{4.40}, one gets
\[ \label{4.43}
\begin{split}
\lim_{n\to\infty}  \|(-\Delta)^{s/2}u_n\|^2_2 = \|(-\Delta)^{s/2}u\|^2_2.
\end{split}
\]
Therefore, $u_n\to u$ in $H^s_{r}(\mathbb{R}^3)$ and $\|u\|_2=c$.
 This completes the proof.
\end{proof}

For $\epsilon >0$, we introduce the set
$$
I^{-\epsilon}_{\mu, \tau}=\{u\in  S_{r}(c):I_{\mu, \tau}(u)\leq -\epsilon\}
\subset H^s_{r}(\mathbb{R}^3).
$$
Because $I_{\mu, \tau}$ is continuous and even on $H^s_{r}(\mathbb{R}^3)$, 
$I^{-\epsilon}_{\mu, \tau}$ is closed and symmetric.

\begin{lemma}\label{lem4.4}
For any fixed $k\in \mathbb{N}$, there exist $\epsilon_k=\epsilon(k)>0$ and 
$\mu:=\mu(k)>0$ such that, for any $0<\epsilon\leq\epsilon_k$ and $\mu \geq \mu_k$, 
one has that $\gamma(I^{-\epsilon}_{\mu, \tau})\geq k$.
\end{lemma}

The proof of Lemma \ref{lem4.4} is similar to \cite[Lemma 3.2]{A5878}, 
so we omit it here.
We define the set
$$
\Sigma_k:=\{\Omega \subset   S_{r}(c): \Omega \text{ is  closed and symmetric, }
\gamma(\Omega)\geq k\},
$$ 
and by Lemma \ref{lem2.161}-(ii), we know that
$$
d_k:=\inf_{\Omega\in \Sigma_k}\sup_{u\in \Omega}I_{\mu, \tau}>-\infty
$$
for all $k\in \mathbb{N}$. To prove Theorem \ref{thm1.1}, we introduce the critical 
value, we define
$$
K_d=\{u\in   S_{r}(c): I'_{\mu, \tau}(u)=0, I_{\mu, \tau}(u)=d\}.
$$
Then, we can derive the following conclusion.

\begin{lemma}\label{lem4.5}
If $d=d_k=d_{k+1}=\cdot\cdot\cdot=d_{k+l}$, 
\[
c\in\Big(0, \min\Big\{\big(\frac{\beta}{\mu}\big)^{\frac{1}{p(1-\delta_{p,s})}}, 
\big(\frac{\Lambda}{D\mu}\big)^{\frac{1}{p(1-\delta_{p,s})}}\Big\}\Big),
\]
$\mu> \mu^*_k=\max\{\mu^*_1, \mu_k\}$, then one has $\gamma(K_d)\geq \ell+1$. 
Especially, $I_{\mu, \tau}(u)$ admits at least $\ell+1$ nontrivial critical points.
\end{lemma}

\begin{proof} For $\epsilon>0$, it is easy to check that 
$I^{-\epsilon}_{\mu, \tau}\in \Sigma$. For any fixed $k\in\mathbb{N}$, by Lemma 
\ref{lem4.4}, there exists $\epsilon_k:=\epsilon(k)>0$ and $\mu_k:=\mu(k)>0$ such that, 
if $0<\epsilon<\epsilon_k$ and $\mu\geq\mu_k$, we have
 $\gamma(I^{-\epsilon_k}_{\mu, \tau})\geq k$. 
 Thus, $I^{-\epsilon_k}_{\mu, \tau}\in \Sigma_k$, and moreover,
$$
d_k\leq \sup_{u\in I^{-\epsilon_k}_{\mu, \tau} }I_{\mu, \tau}(u)=-\epsilon_k<0.
$$
Assume that $0>d=d_k=d_{k+1}=\cdot\cdot\cdot=d_{k+l}$, since 
$c\in\Big(0, \min\Big\{\big(\frac{\beta}{\mu}\big)^{\frac{1}{p(1-\delta_{p,s})}}, 
\big(\frac{\Lambda}{D\mu}\big)^{\frac{1}{p(1-\delta_{p,s})}}\Big\}\Big)$, 
 by  Lemma \ref{lem2.161}-(iii), when $\mu>\mu^*_1>0$ large,   
$I_{\mu, \tau}(u)$ satisfies the $(PS)_d$ condition at the level $d<0$. 
So, $K_d$ is a compact set. By  \cite[Theorem 2.1]{L5898}, we know that the 
restricted function $I_{\mu, \tau}|_{S_r(c)}$ possesses at least $\ell+1$ 
nontrivial critical points. The proof is complete.
\end{proof}

\begin{proof}[Proof of Theorem \ref{thm1.1}]
 Let 
\[
c\in\Big(0, \min\Big\{\big(\frac{\beta}{\mu}\big)^{\frac{1}{p(1-\delta_{p,s})}}, 
\big(\frac{\Lambda}{D\mu}\big)^{\frac{1}{p(1-\delta_{p,s})}}\Big\}\Big),
\]
 $\mu\geq \mu^*_k=\max\{\mu^*_1, \mu_k\}$. From Lemma \ref{lem2.161}-(ii), 
 we see that the critical points of $I_{\mu, \tau}$ found in Lemma \ref{lem4.5} 
are the critical points of $I_\mu$, which completes the proof.
\end{proof}

\section{Proof of Theorem 1.2}

From Lemma \ref{lem3.5}, we see that any critical point of $I_\mu|_{S_r(c)}$ belongs 
to $\mathcal{P}_c$. Consequently, the properties of the manifold $\mathcal{P}_c$ have  
relation to the mini-max structure of $I_\mu|_{ S_{r}(c)}$. For $u\in  S_{r}(c)$ and 
$\theta\in \mathbb{R}$, we introduce the transformation:
\[ %\label{5.1}
(\theta \star u)(x):=e^{\frac{3\theta}{2}}u(e^\theta x), \quad
x\in \mathbb{R}^3,\;\theta \in \mathbb{R}.
\]
It is easy to check that the dilations preserve the $L^2$-norm such that
$\theta \star u\in  S_{r}(c)$, by direct calculation, one has
\begin{align*}%\label{5.2}
I(u,\theta)
&= I_\mu((\theta \star u))\\
&:= \frac{a}{2}e^{2s\theta}\int_{\mathbb{R}^3}|(-\Delta)^{s/2}u|^2dx
+\frac{b}{4}e^{4s\theta}\Big(\int_{\mathbb{R}^3}|(-\Delta)^{s/2}u|^2dx\Big)^2\\
&\quad +\frac{1}{4}e^{(3-2s)\theta}\int_{\mathbb{R}^3}\phi^s_u u^2dx-\frac{\mu}{p}e^{\frac{3(p-2)}{2}\theta}\int_{\mathbb{R}^3}|u|^pdx-
\frac{1}{2^*_s}e^{\frac{3(2^*_s-2)}{2}\theta}\int_{\mathbb{R}^3}|u|^{2^*_s}dx.
\end{align*}

\begin{lemma}\label{lem5.1}
Let $u\in  S_{r}(c)$, then
\begin{itemize}
\item[(i)] $\int_{\mathbb{R}^3}|(-\Delta)^{s/2}(\theta \star u)|^2dx\to 0$ and 
$I_\mu((\theta \star u))\to 0$ as $\theta \to -\infty$;

\item[(ii)] $\int_{\mathbb{R}^3}|(-\Delta)^{s/2}(\theta \star u)|^2dx\to +\infty$ and 
$I_\mu((\theta \star u))\to -\infty$ as $\theta \to +\infty$.
\end{itemize}
\end{lemma}

\begin{proof}
 A direct computation shows that
 \begin{gather*} %\label{5.3}
\int_{\mathbb{R}^3}|(-\Delta)^{s/2}(\theta \star u)|^2dx=e^{2s\theta}
\int_{\mathbb{R}^3}|(-\Delta)^{s/2}u|^2dx, \\
\label{5.4}
\int_{\mathbb{R}^3}|(-\Delta)^{s/2}(\theta \star u)|^2dx\to 0\quad
\text{as }\theta \to -\infty, \\
\label{5.41}
\int_{\mathbb{R}^3}|(-\Delta)^{s/2}(\theta \star u)|^2dx\to +\infty\quad
\text{as }\theta \to +\infty.
\end{gather*}
Notice that
\begin{align*}%\label{5.2}
I_\mu((\theta \star u))
&:= \frac{a}{2}e^{2s\theta}\int_{\mathbb{R}^3}|(-\Delta)^{s/2}u|^2dx
+\frac{b}{4}e^{4s\theta}\Big(\int_{\mathbb{R}^3}|(-\Delta)^{s/2}u|^2dx\Big)^2\\
&\quad +\frac{1}{4}e^{(3-2s)\theta}\int_{\mathbb{R}^3}\phi^s_u u^2dx
 -\frac{\mu}{p}e^{\frac{3(p-2)}{2}\theta}\int_{\mathbb{R}^3}|u|^pdx
 -\frac{1}{2^*_s}e^{\frac{3(2^*_s-2)}{2}\theta}\int_{\mathbb{R}^3}|u|^{2^*_s}dx,
\end{align*}
by $\frac{3(2^*_s-2)}{2}>\frac{3(p-2)}{2}>4s>2s>3-2s$, we infer that
$$
 I_\mu((\theta \star u))\to 0 \quad\text{as }\theta \to -\infty,
 \quad I_\mu((\theta \star u))\to -\infty ~as ~\theta \to +\infty.
$$
This completes the proof.
\end{proof}

\begin{lemma}\label{lem5.2} 
There exist $K=K_c>0$ and $\tilde{c}>0$ such that for all $0<c<\tilde{c}$,
\begin{equation}\label{5.3}
0<\sup_{u\in\mathcal{A}_c}I_\mu(u)<\inf_{u\in\mathcal{B}_c}I_\mu(u),
\end{equation}
where $\mathcal{A}_c£º=\{u\in  S_{r}(c):\int_{\mathbb{R}^3}|
(-\Delta)^{s/2}u|^2dx\leq K_c\}$, 
$\mathcal{B}_c=\{u\in S_{r}(c):\int_{\mathbb{R}^3}|(-\Delta)^{s/2}u|^2dx=2 K_c\}$.
\end{lemma}

\begin{proof}
By Lemma {\ref{lem3.2}}, for any $p\in(2, 2^*_s)$, we have
 \begin{equation}\label{5.4}
\|u\|^p_p\leq C{(p,s)}\|(-\Delta)^{s/2}u\|^{\frac{3}{s}(\frac{p}{2}-1)}_2
\|u\|^{\frac{3}{s}(1-\frac{3-2s}{6}p)}_2, \quad \forall u\in H^s(\mathbb{R}^3).
\end{equation}
By the Sobolev inequality \eqref{2.111} and \eqref{5.4}, for $u, v\in  S_{r}(c)$, 
one has
\begin{align*} %\label{5.5}
I_\mu(v)-I_\mu(u)
&= \frac{a}{2}\int_{\mathbb{R}^3}|(-\Delta)^{s/2}v|^2dx
 -\frac{a}{2}\int_{\mathbb{R}^3}|(-\Delta)^{s/2}u|^2dx\\
&\quad +\frac{b}{4}\Big(\int_{\mathbb{R}^3}|(-\Delta)^{s/2}v|^2dx\Big)^2
 -\frac{b}{4}\Big(\int_{\mathbb{R}^3}|(-\Delta)^{s/2}u|^2dx\Big)^2\\
&\quad +\frac{1}{4}\int_{\mathbb{R}^3}\phi^s_v v^2dx-\frac{1}{4}
 \int_{\mathbb{R}^3}\phi^s_u u^2dx+\frac{\mu}{p}\int_{\mathbb{R}^3}|u|^pdx\\
&\quad -\frac{\mu}{p}\int_{\mathbb{R}^3}|v|^pdx+\frac{1}{2^*_s}\int_{\mathbb{R}^3}|u|^{2^*_s}dx-\frac{1}{2^*_s}\int_{\mathbb{R}^3}|v|^{2^*_s}dx\\
&\geq   \frac{a}{2}\Big(\int_{\mathbb{R}^3}|(-\Delta)^{s/2}v|^2dx
 -\int_{\mathbb{R}^3}|(-\Delta)^{s/2}u|^2dx\Big)\\
&\quad +\frac{b}{4}\Big[\Big(\int_{\mathbb{R}^3}|(-\Delta)^{s/2}v|^2dx\Big)^2
 -\Big(\int_{\mathbb{R}^3}|(-\Delta)^{s/2}u|^2dx\Big)^2\Big]\\
&\quad -\frac{1}{4}\Gamma_sC(p_s, s)^{\frac{3+2s}{3}}\|(-\Delta)^{s/2}u\|^{\frac{3-2s}
{s}}_2c^{\frac{6s-3}{s}}\\
&\quad -\frac{\mu}{p}C(p,s)\|(-\Delta)^{s/2}v\|^{p {\delta_{p,s}}}_2 
 c^{\frac{6-p(3-2s)}{2s}}
 -\frac{S^{-\frac{2^*_s}{2}}}{2^*_s}\|(-\Delta)^{s/2}v\|^{2^*_s}_2.
\end{align*}
Let $\|(-\Delta)^{s/2}u\|^2_2\leq K_c$ and $\|(-\Delta)^{s/2}v\|^2_2=2K_c$, 
here $K_c$ will be determined later. Set
$$
\tilde{c}=\Big(\frac{K_c^{\frac{4s-3}{2s}}}{4\Gamma_s C(p_s, s)^{\frac{3+2s}{3}}}a
 \Big)^{\frac{s}{6s-3}},
$$
by a direct computation, we obtain
\begin{equation}\label{5.6}
\begin{split}
&I_\mu(v)-I_\mu(u) \\
&\geq \frac{a}{2}K_c+\frac{3b}{4}K^2_c-\frac{1}{4}\Gamma_s
  C(p_s, s)^{\frac{3+2s}{3}} K^{\frac{3-2s}{2s}}_c 
  \Big(\frac{K_c^{\frac{4s-3}{2s}}}{4\Gamma_s C(p_s, s)^{\frac{3+2s}{3}}}a
   \Big)^{{\frac{s}{6s-3}}\cdot\frac{6s-3}{s}}\\
&\quad -\frac{\mu}{p}C(p,s)(2K_c)^{\frac{p {\delta_{p,s}}}{2}} 
 \Big(\frac{K_c^{\frac{4s-3}{2s}}}{4\Gamma_s C(p_s, s)^{\frac{3+2s}{3}}}a
 \Big)^{\frac{s}{6s-3}\cdot{\frac{6-p(3-2s)}{2s}}}
 -\frac{S^{-\frac{2^*_s}{2}}}{2^*_s}(2K_c)^{\frac{2^*_s}{2}}\\
&\geq  \frac{a}{2}K_c-\frac{a}{16}K_c-\frac{\mu}{p}2^{\frac{p {\delta_{p,s}}}{2}}
 C(p,s)\Big(\frac{1}{4\Gamma_s C(p_s, s)^{\frac{3+2s}{3}}}a
 \Big)^{{\frac{6-p(3-2s)}{2(6s-3)}}}K^{\frac{(4s-3)(6-p(3-2s))}{4s(6s-3)}}_c 
  K^{\frac{3(p-2)}{4s}}_c\\
&\quad -\frac{S^{-\frac{2^*_s}{2}}}{2^*_s}2^{\frac{2^*_s}{2}}(K_c)^{\frac{2^*_s}{2}}\\
&= \frac{7}{16}aK_c-\frac{\mu 2^{\frac{3(p-2)}{4s}}C(p,s)a^{{\frac{6-p(3-2s)}{2(6s-3)}}}}{p(4\Gamma_s C(p_s, s)^{\frac{3+2s}{3}})^{\frac{6-p(3-2s)}{2(6s-3)}}}K^{\gamma_1}_c K_c-\frac{2^{\frac{2^*_s}{2}}}{2^*_sS^{\frac{2^*_s}{2}}}K^{\frac{2^*_s-2}{2}}_c K_c\\
&\geq   \frac{5}{16}aK_c>0,
\end{split}
\end{equation}
where 
\[
\gamma_1=\frac{[6-p(3-2s)][4s-3]+[3(p-2)-4s][6s-3]}{4s(6s-3)}.
\] 
If we take
$$
K_c=\min\Big\{ \Big(\frac{{p(4\Gamma_s C(p_s, s)^{\frac{3+2s}{3}}) }
 ^{\frac{6-p(3-2s)}{2(6s-3)}}}{16 \mu 2^{\frac{3(p-2)}{4s}}C(p,s)
 a^\frac{6-p(3-2s)}{2(6s-3)}} a\Big)^{\gamma_2}, 
 \Big(\frac{2^*_s S^{\frac{2^*_s}{2}}}{2^{\frac{2^*_s}{2}}16} a
  \Big)^{\frac{2}{2^*_s-2}}\Big\},
$$
with 
\[
\gamma_2=\frac{4s(6s-3)}{[6-p(3-2s)][4s-3]+[3(p-2)-4s][6s-3]},
\]
then, by \eqref{5.6} we deduce that \eqref{5.3} holds. The proof is complete.
\end{proof}

From Lemma \ref{lem5.2}  we can deduce the following conclusion.

\begin{corollary}\label{cor5.1}  
Let $K_c$ and $\tilde{c}$ be given in Lemma \ref{lem5.2}, and $u\in  S_{r}(c)$ with 
$\|(-\Delta)^{s/2}u\|^2_2 \leq K_c$, then $I_\mu(u)>0$.
\end{corollary}

\begin{proof}
A direct computation shows that
\begin{align*} %\label{5.7}
I_\mu(u)
&\geq  \frac{a}{2}\|(-\Delta)^{s/2}u\|^2_2+\frac{b}{4}\|(-\Delta)^{s/2}u\|^4_2
 -\frac{\mu}{p}C(p,s)
c^{\frac{6-p(3-2s)}{2s}}\|(-\Delta)^{s/2}u\|^{\frac{3(p-2)}{2s}}_2\\
&\quad -\frac{S^{-\frac{2^*_s}{2}}}{2^*_s}\|(-\Delta)^{s/2}u\|^{2^*_s}_2>0,
\end{align*}
if $\|(-\Delta)^{s/2}u\|^2_2 \leq K_c$, and the conclusion holds.
\end{proof}

Next, we study the characterizations of the mountain pass levels for 
$I(u,\theta)$ and $I_\mu(u)$. We denote the closed set
$I^d_\mu:=\{u\in  S_{r}(c):I_\mu(u)\leq d\}$.

\begin{proposition}\label{pro5.1} 
Assuming that  $\frac{8s}{3}+2<p< 2^*_s$,  we define
$$
\tilde{c}_\mu (c):=\inf_{\tilde{\gamma}\in \tilde{\Gamma}}\max_{t\in[0,1]}
I(\tilde{\gamma}(t)),~~~{c}_\mu (c):=\inf_{{\gamma}\in{\Gamma}}
\max_{t\in[0,1]}I_{\mu}({\gamma}(t)),
$$
where
\begin{gather*}
\tilde{\Gamma}_c=\{\tilde{\gamma}\in C([0,1],  S_{r}(c)\times \mathbb{R}): 
\tilde{\gamma}(0)\in (\mathcal{A}_c, 0), \tilde{\gamma}(1)\in (I^0_\mu, 0)\},\\
{\Gamma}_c=\{{\gamma}\in C([0,1],  S_{r}(c)): {\gamma}(0)\in \mathcal{A}_c,
 {\gamma}(1)\in I^0_\mu\}.
\end{gather*}
Then we have 
$\tilde{c}_\mu (c)={c}_\mu (c)>0$.
\end{proposition}

\begin{proof}
 On the one hand, for any $\tilde{\gamma}\in \tilde{\Gamma}_c$, we can write  it into 
$$
\tilde{\gamma}(t)=(\tilde{\gamma}_1(t), \tilde{\gamma}_2(t))\in S_r(c)\times \mathbb{R}.
$$
We set ${\gamma}(t)=\tilde{\gamma}_2(t)\star \tilde{\gamma}_1(t)$, then 
${\gamma}\in {\Gamma}_c$, and
$$
 \max_{t\in[0,1]} I(\tilde{\gamma}(t))=\max_{t\in[0,1]} I_\mu(\tilde{\gamma}_2(t)
 \star \tilde{\gamma}_1(t))=\max_{t\in[0,1] } I_\mu({\gamma}(t)),
$$
which implies $\tilde{c}_\mu (c)\geq {c}_\mu (c)>0$, using Corollary \ref{cor5.1}. 
On the other hand, for any ${\gamma}\in {\Gamma}_c$, if we set
$\tilde{\gamma}(t)= ({\gamma}(t), 0)$, then we obtain 
$\tilde{\gamma}\in \tilde{\Gamma}_c$ and
$$
\max_{t\in[0,1]} I(\tilde{\gamma}(t))
=\max_{t\in[0,1] } I_\mu({\gamma}(t)).
$$
This infers that  $\tilde{c}_\mu (c)\leq {c}_\mu (c)$. So,  
$\tilde{c}_\mu (c)= {c}_\mu (c)>0$.
\end{proof}

Next, we show the existence of the $(PS)_{{c}_\mu (c)}$-sequence for $I(u, \theta)$ 
on $ S_{r}(c)\times \mathbb{R}\subset \mathbb{H}$. 
It is obtained by a standard argument using Ekeland's variational principle and
constructing pseudo-gradient flow, see \cite[Proposition 2.2]{J7888}.

\begin{lemma}\label{lem5.4} 
Let $\{\tilde{h}_n\}\subset \tilde{\Gamma}_c $ satisfy that
$$
\max_{t\in[0,1]}I(\tilde{h}_n(t))\leq \tilde{c}_\mu (c)+\frac{1}{n},
$$
then there exists a sequence $\{(v_n, \theta_n)\}\subset S_{r}(c)\times \mathbb{R}$ 
such that
\begin{itemize}
\item[(i)] $I{(v_n, \theta_n)}\in [\tilde{c}_\mu (c)-\frac{1}{n}, 
 \tilde{c}_\mu (c)+\frac{1}{n}]$,
 
\item[(ii)] $\min_{t\in[0,1]}\|(v_n, \theta_n)-\tilde{h}_n(t)\|_{\mathbb{H}}
\leq \frac{1}{\sqrt{n}}$; and

\item[(iii)] $\|(I|_{S_{r}(c)\times \mathbb{R}})'(v_n, \theta_n)\|
 \leq \frac{2}{\sqrt{n}}$, that is,
$$
|\langle I'(v_n, \theta_n), z\rangle|_{\mathbb{H}^{-1}\times \mathbb{H}}
\leq \frac{2}{\sqrt{n}} \|z\|_{\mathbb{H}},
$$
for all
$$ 
{z\in {\tilde{T}_{(v_n, \theta_n)}}}=\{(z_1, z_2)\in \mathbb{H}: 
{\langle v_n, z_1\rangle}_{L^2}=0\}.
$$
\end{itemize}
\end{lemma}


It follows from the above proposition, we can obtain a special 
$(PS)_{{c}_\mu (c)}$-sequence for $I_\mu(u)$ on $S_{r}(c)\subset H^s_r(\mathbb{R}^3)$.

\begin{lemma}\label{lem5.5} 
Under the  assumption $2+\frac{8s}{3}<p<2^*_s$, there exists a sequence
$\{u_n\} \subset S_{r}(c)$ such that
\begin{itemize}
\item[(1)] $I_\mu(u_n)\to c_\mu(c) ~as~ n\to \infty$;

\item[(2)] $P_\mu(u_n)\to 0 ~as~ n\to \infty$;

\item[(3)] $(I_\mu|_{S_{r}(c)})'(u_n)\to 0$  as $n\to \infty$, i.e., 
$\langle I'_\mu(u_n), z\rangle_{{H^{-s}_r} \times H^s_r} \to 0$, 
uniformly for all $z$ satisfying
$$
\|z\|_{H^s_r}\leq 1,\quad \text{where } z\in T_{u_n} 
:=\{z\in {H^s_r(\mathbb{R}^3)}:~\langle u_n,z \rangle_{L^2}=0\}.
$$
\end{itemize}
\end{lemma}

\begin{proof}
Let $\{h_n\}\subset \Gamma_c$ satisfy 
\begin{equation}\label{5.8}
\max_{t\in[0.1]}I_\mu(h_n(t))\leq {c}_\mu(c)+\frac{1}{n},
\end{equation}
we define $\tilde{h}_n(t)=(h_n(t) , 0), \forall t\in[0,1]$. 
It is easy to see that $\tilde{h}_n \in \tilde{\Gamma}_c$ and 
$I_\mu(h_n(t))=I(\tilde{h}_n(t))$. By Proposition \ref{pro5.1}, we have 
$\tilde{c}_\mu(c)={c}_\mu(c)$, then it follows from \eqref{5.8} that
\[  %\label{5.9}
\max_{t\in[0.1]}I(\tilde{h}_n(t))\leq \tilde{c}_\mu(c)+\frac{1}{n}.
\]
It follows from Lemma \ref{lem5.4} that, there exists a sequence
$\{(v_n, \theta_n)\} \subset S_{r}(c)\times \mathbb{R}$ such that as $n\to \infty$,
one has
\begin{gather}\label{5.9}
I(v_n, \theta_n)\to c_\mu(c), ~\theta_n\to 0, \\
\label{5.10}
(I|_{S_{r}(c)\times{\mathbb{R}}})'(v_n, \theta_n)\to 0.
\end{gather}
Set $u_n=\theta_n \star  v_n$. Then, $I_\mu(u_n)=I(v_n, \theta_n)$,
and by \eqref{5.9}, item (1) holds.  To prove conclusion (2), we utilize
\begin{align*}% \label{5.11}
\partial_\theta I(v_n, \theta_n)
&= ase^{2s\theta_n}\|(-\Delta)^{s/2}v_n\|^2_2+bse^{4s\theta_n}\|(-\Delta)^{s/2}v_n\|^4_2+\frac{3-2s}{4}e^{(3-2s)\theta_n}
\int_{\mathbb{R}^3}\phi^s_{v_n}{v^2_n}dx\\
&\quad -\frac{\mu}{p}\frac{3(p-2)}{2}e^{\frac{3(p-2)}{2}\theta n}\int_{\mathbb{R}^3}|v_n|^pdx-\frac{3(2^*_s-2)}{2 2^*_s}e^{\frac{3(2^*_s-2)}{2}\theta n}\int_{\mathbb{R}^3}|v_n|^{2^*_s}dx\\
&= s(a+b\|(-\Delta)^{s/2}u_n\|^2_2)\|(-\Delta)^{s/2}u_n\|^2_2+\frac{3-2s}{4}\int_{\mathbb{R}^3}\phi^s_{u_n}{u^2_n}dx\\
&\quad -\frac{3(p-2)}{2p}\mu \int_{\mathbb{R}^3}|u_n|^pdx-s\int_{\mathbb{R}^3}|u_n|^{2^*_s}dx\\
&= P_\mu(u_n),
\end{align*}
which implies item (2) by \eqref{5.10}. To show item (3), we set $z_n\in T_{u_n}$.
Then,
\begin{align*}%\label{5.12}
\langle I'_\mu(u_n),z_n\rangle
&= a\int_{\mathbb{R}^3}\int_{\mathbb{R}^3}
 \frac{(u_n(x)-u_n(y))(z_n(x)-z_n(y))}{|x-y|^{3+2s}}\,dx\,dy\\
&\quad +b\int_{\mathbb{R}^3}\int_{\mathbb{R}^3}\frac{(u_n(x)
 -u_n(y))^2}{|x-y|^{3+2s}}\,dx\,dy
 \int_{\mathbb{R}^3}\int_{\mathbb{R}^3}
 \frac{(u_n(x)-u_n(y))(z_n(x)-z_n(y))}{|x-y|^{3+2s}}\,dx\,dy  \\
&\quad +\int_{\mathbb{R}^3}\phi^s_{u_n}{u_n}z_ndx
 -\mu\int_{\mathbb{R}^3}|u_n|^{p-2}u_n z_ndx-\int_{\mathbb{R}^3}|u_n|^{2^*_s-2}u_n z_ndx\\
&=  a\int_{\mathbb{R}^3}\int_{\mathbb{R}^3}\frac{(e^{\frac{3\theta_n}{2}}
 v_n(e^{\theta_n} x)-e^{\frac{3\theta_n}{2}}v_n(e^{\theta_n} y))(z_n(x)
 -z_n(y))}{|x-y|^{3+2s}}\,dx\,dy\\
&\quad +b\int_{\mathbb{R}^3}\int_{\mathbb{R}^3}
 \frac{(e^{\frac{3\theta_n}{2}}v_n(e^{\theta_n} x)-e^{\frac{3\theta_n}{2}}
 v_n(e^{\theta_n} y))^2}{|x-y|^{3+2s}}\,dx\,dy \\
 &\quad\times \int_{\mathbb{R}^3}
 \int_{\mathbb{R}^3}\frac{(e^{\frac{3\theta_n}{2}}v_n(e^{\theta_n } x)
 -e^{\frac{3\theta_n}{2}}v_n(e^{\theta_n} y))(z_n(x)-z_n(y))}{|x-y|^{3+2s}}\,dx\,dy\\
&\quad +e^{\frac{3-4s}{2}\theta_n}\int_{\mathbb{R}^3}\phi^s_{v_n}{v_n}
 z_n(e^{-\theta_n}x)dx-\mu e^{\frac{3(p-3)}{2}\theta_n}\int_{\mathbb{R}^3}
 |v_n|^{p-2}v_n z_n(e^{-\theta_n}x) dx\\
&\quad -e^{\frac{3(2^*_s-3)}{2}\theta_n}\int_{\mathbb{R}^3}|v_n|^{2^*_s-2}
 v_n z_n(e^{-\theta_n}x) dx\\
&= e^{2s\theta_n}a\int_{\mathbb{R}^3}\int_{\mathbb{R}^3}\frac{(v_n(x)-v_n(y))
 e^{-\frac{3}{2}\theta_n}(z_n(e^{-\theta_n} x)-z_n(e^{-\theta_n}y))}
  {|x-y|^{3+2s}}\,dx\,dy\\
&\quad +e^{4s\theta_n}b\int_{\mathbb{R}^3}\int_{\mathbb{R}^3}
 \frac{(v_n( x)-v_n( y))^2}{|x-y|^{3+2s}}\,dx\,dy \\
 &\quad\times  \int_{\mathbb{R}^3}\int_{\mathbb{R}^3}\frac{(v_n( x)-v_n( y))
 e^{-\frac{3}{2}\theta_n}(z_n(e^{-\theta_n} x)
 -z_n(e^{-\theta_n}y))}{|x-y|^{3+2s}}\,dx\,dy  \\
&\quad +e^{\frac{3-4s}{2}\theta_n}\int_{\mathbb{R}^3}\phi^s_{v_n}{v_n}z_n
 (e^{-\theta_n}x)dx-\mu e^{\frac{3(p-3)}{2}\theta_n}
 \int_{\mathbb{R}^3}|v_n|^{p-2}v_n z_n(e^{-\theta_n}x) dx\\
&\quad -e^{\frac{3(2^*_s-3)}{2}\theta_n}\int_{\mathbb{R}^3}|v_n|^{2^*_s-2}
v_n z_n(e^{-\theta_n}x) dx.
\end{align*}
Denoting $\tilde{z}_n(x)=e^{-\frac{3\theta_n}{2}}z_n(e^{-\theta_n}x)$,
 we obtain
$$
\langle I'_\mu(u_n),z_n\rangle_{{H^{-s}_r}\times H^s_r}
=\langle I'(v_n, \theta_n), (\tilde{z}_n, 0) \rangle_{\mathbb{H}^{-1}\times \mathbb{H}}.
$$
It is easy to check that
\[  %\label{5.12}
 \langle v_n, \tilde{z}_n\rangle_{L^2}=\int_{\mathbb{R}^3} v_n(x) e^{-\frac{3\theta_n}{2}}z_n(e^{-\theta_n}x)dx
 =\int_{\mathbb{R}^3} v_n(e^{\theta_n}x) e^{\frac{3\theta_n}{2}}z_n(x)dx
 =\int_{\mathbb{R}^3}u_n(x)z_n(x)dx=0.
\]
Therefore, $(\tilde{z}_n, 0)\in \tilde{T}_{(v_n, \theta_n)}$.
On the other hand,
$$
\|(\tilde{z}_n,0)\|^2_{\mathbb{H}}=\|\tilde{z}_n\|^2_{{H^s_r}}=\|z_n\|^2_2
+e^{-2s\theta_n}\|z_n\|^2_{D^{s,2}}\leq C\|z_n\|^2_{H^s_r},
$$
where the last inequality follows by $\theta_n\to 0$. Consequently, we conclude
item (3). The proof is complete.
\end{proof}


\begin{lemma}\label{lem5.6} 
The $(PS)$ sequence $\{u_n\} \subset S_{r}(c)$ for $I_\mu(u)$ with the level 
$c_\mu(c)$ mentioned in Lemma \ref{lem5.5} is bounded in $H^s_r(\mathbb{R}^3)$.
\end{lemma}

\begin{proof}  
From Lemma  \ref{lem5.5} (1), we see that $I_\mu(u)$ is bounded. 
In fact, by $P_\mu(u_n)\to 0$ as $n\to \infty$, one obtains
$$
|(1+2s)I_\mu(u_n)+P_\mu(u_n)|\leq 3c_\mu(c),
$$
which implies that
\begin{equation}\label{5.13}
\begin{split}
-3c_\mu(c)&\leq \frac{1+4s}{2}a\|(-\Delta)^{s/2}u_n\|^2_2+\frac{1+6s}{4}b
\|(-\Delta)^{s/2}u_n\|^4_2+\int_{\mathbb{R}^3}\phi^s_{u_n}u^2_ndx\\
&\quad -\mu(\frac{1+2s}{p}+s\delta_{p,s})\int_{\mathbb{R}^3}|u_n|^pdx
 -(\frac{1+2s}{2^*_s}+s)\int_{\mathbb{R}^3}|u_n|^{2^*_s}dx.
\end{split}
\end{equation}
In view of the boundedness of $I_\mu(u)$, we have
\begin{equation}\label{5.14}
\begin{split}
&a\|(-\Delta)^{s/2}u_n\|^2_2+\frac{b}{2}\|(-\Delta)^{s/2}u_n\|^4_2
 +\frac{1}{2}\int_{\mathbb{R}^3}\phi^s_{u_n}u^2_ndx\\
&\leq 6c_\mu(c)+\frac{2\mu}{p}\int_{\mathbb{R}^3}|u_n|^pdx
 +\frac{2}{2^*_s}\int_{\mathbb{R}^3}|u_n|^{2^*_s}dx.
\end{split}
\end{equation}
By \eqref{5.13} and \eqref{5.14}, we obtain
\begin{align*} %\label{5.15}
\mu \frac{(p\delta_{p,s}-4)s}{p}\int_{\mathbb{R}^3}|u_n|^pdx
+\frac{(2^*_s-4)s}{2^*_s}\int_{\mathbb{R}^3}|u_n|^{2^*_s}dx
+\frac{(6s-3)}{4}\int_{\mathbb{R}^3}\phi^s_{u_n}u^2_ndx\leq 3(2+6s)c_\mu(c).
\end{align*}
Note that $s\in(\frac{3}{4}, 1)$, $p>\frac{8s}{3}+2$, we have that 
$p\delta_{p,s}-4>0$, and so
$$
\int_{\mathbb{R}^3}|u_n|^pdx, \int_{\mathbb{R}^3}|u_n|^{2^*_s}dx 
\quad\text{and}\quad
\int_{\mathbb{R}^3}\phi^s_{u_n}u^2_ndx 
$$
are all bounded. Thus, $\|(-\Delta)^{s/2}u_n\|_2\leq R_3$ for some $R_3>0$ 
independently on $n\in \mathbb{N}$. Since $\{u_n\}\subset S_{r}(c)$, we see that 
$\{u_n\}$ is bounded in $H^s_r(\mathbb{R}^3)$.  
This completes the proof.
\end{proof}

Now, we set the functional $\Phi:H^s_r(\mathbb{R}^3)\to \mathbb{R}$ as
$$
\Phi(u)=\frac{1}{2}\int_{\mathbb{R}^3} |u|^2dx,
$$
then $S_{r}(c)=\Phi^{-1}(\frac{c^2}{2})$. As a result, it can be derived from 
\cite[Proposition 5.12]{W6888} that there is a sequence 
$\{\lambda_n\}\subset \mathbb{R}$ such that
$$
I'_\mu(u_n)-\lambda_n \Phi'(u_n)\to 0,~in ~H^{-s}_r(\mathbb{R}^3) ~as~n\to \infty.
$$
That is, in $H^{-s}_r(\mathbb{R}^3)$, we have
\begin{equation}\label{5.16}
\left(a+b\|(-\Delta)^{s/2}u_n\|^2_2\right)(-\Delta)^s u_n+\phi^s_{u_n}u_n
-\mu |u_n|^{p-2}u_n-|u_n|^{2^*_s-2}u_n
=\lambda_n u_n+o_n(1).
\end{equation}
Therefore, for any $\varphi \in H^s_r(\mathbb{R}^3)$, one has
\begin{equation}\label{5.166}
\begin{split}
&\left(a+b\|(-\Delta)^{s/2}u_n\|^2_2\right)
\int_{\mathbb{R}^3}(-\Delta)^{s/2} u_n(-\Delta)^{s/2} \varphi dx
+\int_{\mathbb{R}^3}\phi^s_{u_n}u_n\varphi dx \\
&-\int_{\mathbb{R}^3}\mu |u_n|^{p-2}u_n\varphi dx
 -\int_{\mathbb{R}^3}|u_n|^{2^*_s-2}u_n\varphi dx \\
&=\lambda_n \int_{\mathbb{R}^3}u_n\varphi dx +o_n(1).
\end{split}
\end{equation}

Next, we study the asymptotical behavior of the mountain pass level value 
$c_\mu(c)$ as $\mu \to +\infty$, and the properties of  the $(PS)_{c_\mu(c)}$-sequence 
$\{u_n\}\subset S_{r}(c)$ as $n\to +\infty$.

\begin{lemma}\label{lem5.7} 
The limit $\lim_{\mu\to +\infty}c_\mu(c)=0 $ holds.
\end{lemma}

\begin{proof} 
Recall Lemma {\ref{lem5.1}} and Corollary {\ref{cor5.1}}, we see that for fixed 
$u_0\in S_{r}(c)$, there exists two constants
$\theta_1, \theta_2$ satisfying $\theta_1<0< \theta_2$ such that 
$u_1=\theta_1 \star u_0 \in \mathcal{A}_c$ and 
$I_\mu(u_2)=I_\mu((\theta_2 \star u_0 ))<0$. Then, we can define a path
$$
\eta_0: t \in [0,1]\to ((1-t)\theta_1+t \theta_2)\star u_0 \in \Gamma_c.
$$
Therefore,
\begin{align*}
0<c_\mu(c)
&\leq  \max_{t\in[0,1]} I_\mu(\eta_0(t))\\
&\leq  \max_{r\geq 0} \Big\{\frac{a}{2}r^{2s}\|(-\Delta)^{s/2}u_0\|^2_2+\frac{b}{4}r^{4s}\|(-\Delta)^{s/2}u_0\|^4_2\\
&\quad +\frac{1}{4}r^{3-2s}\int_{\mathbb{R}^3}\phi^s_{u_0}u^2_0dx
 -\frac{\mu}{p}r^{\frac{3(p-2)}{2}}\int_{\mathbb{R}^3}|u_0|^pdx\Big\}\\
&:=\max_{r\geq 0}g(r).
\end{align*}
Note that $\frac{3(p-2)}{2}>4s>2s>3-2s$, we have that $\lim_{r\to 0^+}g(r)=0^+$, 
$\lim_{r\to +\infty}g(r)=-\infty$. Then, there exists a unique maximum point $r_0>0$ 
such that $\max_{r\geq 0}g(r)=g(r_0)>0$. Hence, we distinguish two cases: $r_0\geq1$ 
and $0 \leq r_0<1$.

If $r_0\geq1$, then by $s\in(\frac{3}{4}, 1)$, we have
\begin{align*}
\max_{t\in[0,1]} I_\mu(\eta_0(t))&\leq  g(r_0)\\
&\leq  \Big\{\frac{a}{2}r^{4s}_0\|(-\Delta)^{s/2}u_0\|^2_2
 +\frac{b}{4}r^{4s}_0\|(-\Delta)^{s/2}u_0\|^4_2\\
&\quad +\frac{1}{4}r^{4s}_0\int_{\mathbb{R}^3}\phi^s_{u_0}u^2_0dx
 -\frac{\mu}{p}r^{\frac{3(p-2)}{2}}_0\int_{\mathbb{R}^3}|u_0|^pdx\Big\}\\
&\leq  \max_{r\geq 0}\Big\{3\max\Big\{\frac{a}{2}\|(-\Delta)^{s/2}u_0\|^2_2, 
 \frac{b}{4}\|(-\Delta)^{s/2}u_0\|^4_2,
 \frac{1}{4}\int_{\mathbb{R}^3}\phi^s_{u_0}u^2_0dx  \Big\}r^{4s}\\
&\quad -\frac{\mu}{p}r^{\frac{3(p-2)}{2}}\int_{\mathbb{R}^3}|u_0|^pdx\Big\}\\
&= 3m(r_{\rm max})^{4s}-\frac{\mu}{p}n(r_{\rm max})^{\frac{3(p-2)}{2}}\\
&= \frac{m(3p-6-8s)}{p-2}\left[\frac{8psm}{(p-2)\mu n}\right]^{\frac{8s}{3p-6-8s}},
\end{align*}
where $r_{\rm max}=\left[\frac{8psm}{(p-2)\mu n}\right]^{\frac{2}{3p-6-8s}}, 
~m=\max\left\{\frac{a}{2}\|(-\Delta)^{s/2}u_0\|^2_2, \frac{b}{4}\|(-\Delta)^{s/2}u_0\|^4_2,
\frac{1}{4}\int_{\mathbb{R}^3}\phi^s_{u_0}u^2_0dx  \right\},$ $n=\int_{\mathbb{R}^3}|u_0|^pdx.$
Therefore, for $\frac{8s}{3}+2<p<2^*_s$, we have a positive constant $\tilde{C}$ independent of $\mu$ such that
$$c_\mu (c)\leq \tilde{C}{\mu}^{-\frac{8s}{3p-6-8s}}\to 0, ~~as~\mu\to +\infty.$$


If $0 \leq r_0<1$, we infer to
\begin{align*}
\max_{t\in[0,1]} I_\mu(\eta_0(t))&\leq  g(r_0)\\
&\leq  \Big\{\frac{a}{2}r^{3-2s}_0\|(-\Delta)^{s/2}u_0\|^2_2
 +\frac{b}{4}r^{3-2s}_0\|(-\Delta)^{s/2}u_0\|^4_2\\
&\quad +\frac{1}{4}r^{3-2s}_0\int_{\mathbb{R}^3}\phi^s_{u_0}u^2_0dx
 -\frac{\mu}{p}r^{\frac{3(p-2)}{2}}_0\int_{\mathbb{R}^3}|u_0|^pdx\Big\}\\
&\leq  \max_{r\geq 0}\Big\{3\max\Big\{\frac{a}{2}\|(-\Delta)^{s/2}u_0\|^2_2, 
 \frac{b}{4}\|(-\Delta)^{s/2}u_0\|^4_2,
\frac{1}{4}\int_{\mathbb{R}^3}\phi^s_{u_0}u^2_0dx  \Big\}r^{3-2s}\\
&\quad -\frac{\mu}{p}r^{\frac{3(p-2)}{2}}\int_{\mathbb{R}^3}|u_0|^pdx\Big\}\\
&= 3m(r_{\rm max})^{3-2s}-\frac{\mu}{p}n(r_{\rm max})^{\frac{3(p-2)}{2}}\\
&= \frac{m(3p+4s-12)}{p-2}\Big[\frac{2pm(3-2s)}{(p-2)\mu n}\Big]^{\frac{6-4s}{3p+4s-12}},
\end{align*}
where $r_{\rm max}=[\frac{2pm(3-2s)}{(p-2)\mu n}]^{\frac{2}{3p+4s-12}}$.
Therefore, for $2+\frac{8s}{3}<p<2^*_s$, and $s\in{(\frac{3}{4}, 1)}$, we can deduce 
that $3p+4s-12>0$, then there exists a positive constant $\bar{C}$ independent of
 $\mu$ such that
$$
c_\mu (c)\leq \bar{C}{\mu}^{-\frac{6-4s}{3p+4s-12}}\to 0, ~~as~\mu\to +\infty.
$$
This completes the proof.
\end{proof}

\begin{lemma}\label{lem5.8} 
There exists  a constant $C=C(p,s)>0$ such that
\begin{gather*}
\limsup_{n\to \infty} \int_{\mathbb{R}^3}F(u_n)dx\leq C c_\mu(c), \quad
\limsup_{n\to \infty} \int_{\mathbb{R}^3}f(u_n){u_n}dx\leq C c_\mu(c), \\
\limsup_{n\to \infty} \int_{\mathbb{R}^3}\phi^s_{u_n}u^2_ndx\leq C c_\mu(c),\quad 
\limsup_{n\to \infty} \int_{\mathbb{R}^3}|(-\Delta)^{s/2}u_n|^2dx\leq C c_\mu(c),\\
\limsup_{n\to \infty} (\int_{\mathbb{R}^3}|(-\Delta)^{s/2}u_n|^2dx)^2\leq C c_\mu(c),
\end{gather*}
where $f(u)=\mu|u|^{p-2}u+|u|^{2^*_s-2}u$, $F(u)=\int^u_0 f(s)ds$.
\end{lemma}

\begin{proof} 
Since $I_\mu(u_n)\to c_\mu(c)$ and $P_\mu(u_n)\to 0$ as $n \to \infty$, one obtains
\begin{equation}\label{5.17}
\begin{split}
&3c_\mu(c)+o_n {(1)}=3I_\mu(u_n)+P_\mu(u_n)\\
&= \frac{3+2s}{2}a\int_{\mathbb{R}^3}|(-\Delta)^{s/2}u_n|^2dx
+\frac{3+4s}{4}b\Big(\int_{\mathbb{R}^3}|(-\Delta)^{s/2}u_n|^2dx\Big)^2\\
&\quad +\frac{3-s}{2}\int_{\mathbb{R}^3}\phi^s_{u_n} u^2_ndx
 -\frac{3}{2}\int_{\mathbb{R}^3}f(u_n)u_ndx\\
&\leq  \frac{3+4s}{2}a\int_{\mathbb{R}^3}|(-\Delta)^{s/2}u|^2dx
+\frac{3+4s}{4}b\Big(\int_{\mathbb{R}^3}|(-\Delta)^{s/2}u|^2dx\Big)^2\\
&\quad +\frac{3-s}{2}\int_{\mathbb{R}^3}\phi^s_{u_n} u^2_ndx-\frac{3}{2}\int_{\mathbb{R}^3}f(u_n)u_ndx\\
&= (3+4s)\Big(-\frac{1}{4}\int_{\mathbb{R}^3}\phi^s_{u_n} u^2_ndx
 +\int_{\mathbb{R}^3}F(u_n)dx+c_\mu(c)+o_n(1)\Big)\\
&\quad +\frac{3-s}{2}\int_{\mathbb{R}^3}\phi^s_{u_n} u^2_ndx
 -\frac{3}{2}\int_{\mathbb{R}^3}f(u_n)u_ndx\\
&= (3+4s)\Big(\int_{\mathbb{R}^3}F(u_n)dx+c_\mu(c)+o_n(1)\Big)\\
&\quad -\frac{6s-3}{4}\int_{\mathbb{R}^3}\phi^s_{u_n} u^2_ndx-\frac{3}{2}\int_{\mathbb{R}^3}f(u_n)u_ndx\\
&\leq  (3+4s)\Big(\int_{\mathbb{R}^3}F(u_n)dx+c_\mu(c)+o_n(1)\Big)
 -\frac{3}{2}\int_{\mathbb{R}^3}f(u_n)u_ndx\\
&\leq  (3+4s)(c_\mu(c)+o_n(1))-\frac{3}{2}p\int_{\mathbb{R}^3}F(u_n)dx
 +(3+4s)\int_{\mathbb{R}^3}F(u_n)dx.
\end{split}
\end{equation}
Hence,
$$
4sc_\mu(c)+o_n(1)\geq \frac{3p-6-8s}{2}\int_{\mathbb{R}^3}F(u_n)dx,
$$
which implies that
\begin{equation}\label{5.18}
\limsup_{n\to \infty} \int_{\mathbb{R}^3}F(u_n)dx
\leq \frac{8s}{3p-6-8s}c_\mu(c) 
\leq C c_\mu(c)
\end{equation}
and then
\begin{equation}\label{5.19}
\limsup_{n\to \infty} \int_{\mathbb{R}^3}f(u_n)u_ndx \leq C c_\mu(c).
\end{equation}
Then, from \eqref{5.17}-\eqref{5.19}, one has
\begin{align*} %\label{5.20}
&\limsup_{n\to \infty}\Big\{\frac{3+2s}{2}a
 \int_{\mathbb{R}^3}|(-\Delta)^{s/2}u_n|^2dx
 +\frac{3+4s}{4}b(\int_{\mathbb{R}^3}|(-\Delta)^{s/2}u_n|^2dx)^2
 +\frac{3-s}{2}\int_{\mathbb{R}^3}\phi^s_{u_n} u^2_ndx\Big\}\\
&= \limsup_{n\to \infty}\Big\{\frac{3}{2}\int_{\mathbb{R}^3}f(u_n)u_ndx+3c_\mu(c)\Big\}\\
&\leq  C c_\mu(c).
\end{align*}
Consequently, the proof is complete.
\end{proof}

\begin{lemma}\label{rem5.91} 
Let $\{u_n\} \subset S_r(c)$ be the  $(PS)$ sequence for the constrained functional
 $I_\mu|_{S_r(c)}$ at level 
 $c_\mu(c)\in \left(0, \Big(\frac{1}{4}-\frac{1}{2^*_s}\Big)(aS)^{\frac{3}{2s}}\right)$ 
with $P_\mu(u_n)\to 0$ as $n\to \infty$. Then
\begin{itemize}
\item[(i)] $\{\lambda_n\}$ is bounded in $\mathbb{R}$, and 
$\limsup_{n\to \infty}|\lambda_n|\leq \frac{C}{c^2}c_\mu(c)$ has the  
estimation
$$
\lambda_n=\frac{1}{c^2}\Big[\frac{6s-3}{4s}\int_{\mathbb{R}^3}\phi^s_{u_n} u^2_ndx
+\mu \frac{p(3-2s)-6}{2ps}\int_{\mathbb{R}^3}|u_n|^pdx\Big]+o_n(1).
$$
Moreover, there exists some $\mu^*_2:=\mu^*_2(c)>0$ such that 
$\lim_{n\to +\infty}\lambda_n=\lambda<0$, if $\mu> \mu^*_2$ large;

\item[(ii)] there exist $u\in S_r(c)$ such that, up to a subsequence, 
$u_n\to u$ strongly in $H^s_r(\mathbb{R}^3)$ as  $\mu> \mu^*_2$ large and 
$u$ is a solution of  system \eqref{lab1.1} for some $\lambda<0$.
\end{itemize}
\end{lemma}

\begin{proof}  We split the proof into three steps.
\smallskip

\noindent\textbf{Step 1.} 
We assert that  $u\not\equiv 0$.
From Lemma {\ref{lem5.6}}, we  know that $\{u_n\}$ is a  bounded $(PS)$ sequence 
for $I_\mu$ in $H^s_r(\mathbb{R}^3)$, and by Lemma {\ref{lem2.1311}}, up to a 
subsequence, there exists $u\in H^s_r(\mathbb{R}^3)$ such that $u_n\rightharpoonup u$ 
weakly in $ H^s_r(\mathbb{R}^3)$, $u_n\to u$ strongly in $ L^p(\mathbb{R}^3)$, for 
$p\in (2, 2^*_s)$, $u_n(x)\rightharpoonup u(x)$ a.e.\ on $\mathbb{R}^3$. 
In view of $2+\frac{8s}{3}<p<2^*_s$, and Lemma {\ref{lem2.1311}} and  
Lemma {\ref{lem2.171}}, then
\begin{equation}\label{5.2122}
\lim_{n\to \infty}\int_{\mathbb{R}^3}|u_n|^pdx
=\int_{\mathbb{R}^3}|u|^pdx,~\lim_{n\to \infty}\int_{\mathbb{R}^3}
\phi^s_{u_n}u^2_ndx=\int_{\mathbb{R}^3}\phi^s_{u}u^2dx.
\end{equation}
Suppose by contradiction that, $u\equiv 0$. Then, by \eqref{5.2122} and 
$P_\mu(u_n)=o_n(1)$, we deduce that
\begin{align*}
o_n(1)&
= \Big(a+b\int_{\mathbb{R}^3}|(-\Delta)^{s/2}u_n|^2dx\Big)
 \int_{\mathbb{R}^3}|(-\Delta)^{s/2}u_n|^2dx
 +\frac{3-2s}{4s}\int_{\mathbb{R}^3}\phi^s_u u^2dx\\
&\quad -\mu \delta_{p,s}\int_{\mathbb{R}^3}|u_n|^pdx
 -\int_{\mathbb{R}^3}|u_n|^{2^*_s}dx\\
&= \Big(a+b\int_{\mathbb{R}^3}|(-\Delta)^{s/2}u_n|^2dx\Big)
 \int_{\mathbb{R}^3}|(-\Delta)^{s/2}u_n|^2dx
-\int_{\mathbb{R}^3}|u_n|^{2^*_s}dx+o_n(1).
\end{align*}
Without loss of generality, we assume that
$$
a\int_{\mathbb{R}^3}|(-\Delta)^{s/2}u_n|^2dx
+b\Big(\int_{\mathbb{R}^3}|(-\Delta)^{s/2}u_n|^2dx\Big)^2\to l\geq 0,\quad
\int_{\mathbb{R}^3}|u_n|^{2^*_s}dx\to l,
$$
as $n\to \infty$. If $l=0$, then we can deduce from the expression of $I_\mu(u_n)$ 
that $c_\mu(c)=0$, which is absurd since $c_\mu(c)>0$. Hence, $l>0$. 
By the definition of $S$, we have
$$
S\leq \frac{\int_{\mathbb{R}^3}|(-\Delta)^{s/2}u_n|^2dx}
{\|u_n\|^2_{2^*_s}}\leq \frac{1}{a}
\frac{a\int_{\mathbb{R}^3}|(-\Delta)^{s/2}u_n|^2dx
+b\Big(\int_{\mathbb{R}^3}|(-\Delta)^{s/2}u_n|^2dx\Big)^2}
{\left(\int_{\mathbb{R}^3}|u_n|^{2^*_s}
dx\right)^{2/2^*_s}}
\to \frac{1}{a}l^{\frac{2s}{3}}
$$
as $n\to \infty $. It follows that $l\geq (aS)^{\frac{3}{2s}}$. Consequently,
 by \eqref{5.2122} we have
\begin{align*}
c_\mu(c)
&= \lim_{n\to \infty}I_\mu(u_n)\\
&= \lim_{n\to \infty}\Big\{\frac{a}{2}\int_{\mathbb{R}^3}|(-\Delta)^{s/2}u_n|^2dx
 +\frac{b}{4}\Big(\int_{\mathbb{R}^3}|(-\Delta)^{s/2}u_n|^2dx\Big)^2 \\
&+ \frac{1}{4}\int_{\mathbb{R}^3}\phi^s_{u_n} u^2_ndx
  -\frac{\mu}{p}\int_{\mathbb{R}^3}|u_n|^pdx
  -\frac{1}{2^*_s}\int_{\mathbb{R}^3}|u_n|^{2^*_s}dx\Big\}\\
&\geq \Big(\frac{1}{4}-\frac{1}{2^*_s}\Big) l \\
&\geq \Big(\frac{1}{4}-\frac{1}{2^*_s}\Big)(aS)^{\frac{3}{2s}},
\end{align*}
which contradicts $I_\mu(u_n)\to c_\mu(c)
<\Big(\frac{1}{4}-\frac{1}{2^*_s}\Big)(aS)^{\frac{3}{2s}}$. 
Therefore, $u\not \equiv 0$.
\smallskip

\noindent\textbf{Step 2.} 
We prove that $u_n\to u$ in $L^{2^*_s}(\mathbb{R}^3)$.
Again by Lemma \ref{lem5.6}, we can obtain 
$\{\|(-\Delta)^{s/2}u_n\|_2\}$ is bounded in $\mathbb{R}$, by Prohorov's theorem 
\cite{P9899}, there exist two positive measures, $\zeta$,
 $\omega\in \mathcal{M}(\mathbb{R}^3)$, such that
\[ 
|(-\Delta)^{s/2}u_n|^2\rightharpoonup \omega, |u_n|^{2^*_s}\rightharpoonup\zeta \quad
\text{in } \mathcal{M}(\mathbb{R}^3)
\]
as $n\to \infty$. Then, by Lemma \ref{lem2.131}, either $u_n\to u$ in
$L^{2^*_s}_{loc}(\mathbb{R}^3)$ or there exists a (at most countable) set of distinct
points $\{x_j\}_{j\in J}\subset \mathbb{R}^3$ and positive numbers $\{\zeta_j\}_{j\in J}$ such that
$$
\zeta=|u|^{2^*_s}+\Sigma_{j\in J}\zeta_j \delta_{x_j}.
$$
Moreover, there exist some at most a countable set $J\subset \mathbb{N}$,
a corresponding set of distinct points $\{x_j\}_{j\in J}\subset \mathbb{R}^3$,
and two sets of positive numbers  $\{\zeta_j\}_{j\in J}\subset \mathbb{R}^3$ and
$\{\omega_j\}_{j\in J}\subset \mathbb{R}^3$ such that items
\eqref{lab3.11}-\eqref{lab3.13} hold. Now, assume that $J\neq \emptyset$.

 Similar to the proof in step 1 and step 2 of Lemma \ref{lem2.161},  we  
 can obtain
 \begin{equation}\label{4.1667}
 a\omega(\{x_j\})\leq\zeta_j,  \quad a\omega_\infty\leq\zeta_\infty.
 \end{equation}
Next, we claim that  $\zeta_j=0$ for any $j\in J$ and $\zeta_\infty=0$.

 Suppose by contradiction that, there exists $j_1\in J$ such that $\zeta_{j_1}>0$ 
or $\zeta_\infty>0$.  By  \eqref{4.1667}, Lemma {\ref{lem2.131}} and  
Lemma  \ref{lem2.11}, we obtain
$$
\zeta_{j_1}\geq (aS)^{\frac{3}{2s}} \quad\text{or}\quad
\zeta_{\infty}\geq (aS)^{\frac{3}{2s}}.
$$
If the former case occurs, one has
\begin{align*}
\Big(\frac{1}{4}-\frac{1}{2^*_s}\Big)(aS)^{\frac{3}{2s}}
&>c_\mu(c)
=\lim_{n\to \infty}\Big[I_\mu(u_n)-\frac{1}{4s}P_\mu(u_n)\Big]\\
&=  \lim_{n\to \infty}\big[\frac{1}{4}a\|(-\Delta)^{s/2}u_n\|^2_2
 +\Big(\frac{1}{4}-\frac{3-2s}{16s}\Big)\int_{\mathbb{R}^3}\phi^s_{u_n}u^2_ndx\\
&\quad +\Big(\frac{3(p-2)}{8ps}-\frac{1}{p}\Big)
 \mu \int_{\mathbb{R}^3}|u_n|^pdx
 +\Big(\frac{1}{4}-\frac{1}{2^*_s}\Big)\int_{\mathbb{R}^3}|u_n|^{2^*_s}dx\Big]\\
&=  \lim_{n\to \infty}\Big[\frac{1}{4}a\|(-\Delta)^{s/2}u_n\|^2_2
 +\frac{6s-3}{16s}\int_{\mathbb{R}^3}\phi^s_{u_n}u^2_ndx\\
&\quad +\frac{3(p-2)-8s}{8ps}\mu \int_{\mathbb{R}^3}|u_n|^pdx+\Big(\frac{1}{4}
 -\frac{1}{2^*_s}\Big)\int_{\mathbb{R}^3}|u_n|^{2^*_s}dx\Big]\\
&\geq \lim_{n\to \infty}\Big(\frac{1}{4}-\frac{1}{2^*_s}\Big)
 \int_{\mathbb{R}^3}|u_n|^{2^*_s}dx\\
&\geq \Big(\frac{1}{4}-\frac{1}{2^*_s}\Big)\lim_{\rho\to 0}
 \lim_{n\to \infty}\int_{\mathbb{R}^3}|u_n|^{2^*_s}\varphi_\rho dx\\
&=  \Big(\frac{1}{4}-\frac{1}{2^*_s}\Big)\zeta_{j_1}\\
&\geq\Big(\frac{1}{4}-\frac{1}{2^*_s}\Big) (aS)^{\frac{3}{2s}},
\end{align*}
which is a contradiction. If the latter case happens, we have
\begin{align*}
\Big(\frac{1}{4}-\frac{1}{2^*_s}\Big)(aS)^{\frac{3}{2s}}>c_\mu(c)
&\geq \lim_{n\to \infty}\Big(\frac{1}{4}-\frac{1}{2^*_s}\Big)
 \int_{\mathbb{R}^3}|u_n|^{2^*_s}dx\\
&\geq \Big(\frac{1}{4}-\frac{1}{2^*_s}\Big)\lim_{R\to \infty}
 \lim_{n\to \infty}\int_{\mathbb{R}^3}|u_n|^{2^*_s}\psi_R dx\\
&= \Big(\frac{1}{4}-\frac{1}{2^*_s}\Big)\zeta_\infty \\
&\geq\Big(\frac{1}{4}-\frac{1}{2^*_s}\Big) (aS)^{\frac{3}{2s}}.
\end{align*}
This is also a contradiction. Therefore, $\zeta_j=0$ for any 
$j\in J$ and $\zeta_\infty=0$. As a result, by Lemma \ref{lem2.131}, 
we obtain that $u_n\to u$ in $L^{2^*_s}_{loc}(\mathbb{R}^3)$.
Combining this with   Lemma \ref{lem2.11}, we know that $u_n\to u$ in 
$L^{2^*_s}(\mathbb{R}^3)$.
\smallskip

\noindent\textbf{Step 3.}
We prove that there exists some $\mu^*_2:=\mu^*_2(c)>0$ such that 
$\lim_{n\to +\infty}\lambda_n=\lambda<0$, if $\mu> \mu^*_2$ large.

By \eqref{5.16} and the fact that $u_n\in S_{r}(c)$, one has
\begin{align*}
&a\int_{\mathbb{R}^3}|(-\Delta)^{s/2}u_n|^2dx
 +b\Big(\int_{\mathbb{R}^3}|(-\Delta)^{s/2}u_n|^2dx\Big)^2
+\int_{\mathbb{R}^3}\phi^s_{u_n} u^2_ndx-\int_{\mathbb{R}^3}f(u_n)u_ndx\\
&= \lambda_n \int_{\mathbb{R}^3}|u_n|^2dx+o_n(1)\\
&= \lambda_n c^2+o_n(1).
\end{align*}
It indicates that
\begin{align*}
\lambda_n
&= \frac{1}{ c^2}\Big[a\int_{\mathbb{R}^3}|(-\Delta)^{s/2}u_n|^2dx
+b\Big(\int_{\mathbb{R}^3}|(-\Delta)^{s/2}u_n|^2dx\Big)^2\\
&\quad +\int_{\mathbb{R}^3}\phi^s_{u_n} u^2_ndx
 -\int_{\mathbb{R}^3}f(u_n)u_ndx\Big]+o_n(1).
\end{align*}
By Lemma \ref{lem5.6}, we  know that $\{u_n\}$ is bounded in $H^s_r(\mathbb{R}^3)$, 
and so, $\{\lambda_n\}$ is bounded in $\mathbb{R}$. 
By Lemma {\ref{lem5.8}} we know that 
$\limsup_{n\to \infty}|\lambda_n|\leq \frac{C}{c^2}c_\mu(c)$. 
From this and  $P_\mu(u_n)\to 0$ as $n\to \infty$, we derive that
\begin{align*}
\lambda_n
&= \frac{1}{ c^2}\Big[a\int_{\mathbb{R}^3}|(-\Delta)^{s/2}u_n|^2dx
 +b\Big(\int_{\mathbb{R}^3}|(-\Delta)^{s/2}u_n|^2dx\Big)^2\\
&\quad +\int_{\mathbb{R}^3}\phi^s_{u_n} u^2_ndx
 -\int_{\mathbb{R}^3}f(u_n)u_ndx-\frac{1}{s}P_\mu(u_n)\Big]+o_n(1)\\
&= \frac{1}{c^2}\Big[\frac{6s-3}{4s}\int_{\mathbb{R}^3}\phi^s_{u_n} u^2_ndx
 +\mu \frac{p(3-2s)-6}{2ps}\int_{\mathbb{R}^3}|u_n|^pdx\Big]+o_n(1).
\end{align*}
By \eqref{5.2122} and similar arguments to that of \eqref{4.32}-\eqref{4.35}, 
we see that there exists $\mu^*_2:=\mu^*_2(c)>0$, such that
\begin{equation}\label{5.212}
\begin{split}
\lambda&= \lim_{n\to \infty}\lambda_n\\
&= \lim_{n\to \infty} \frac{1}{c^2} \Big\{\frac{6s-3}{4s}
 \int_{\mathbb{R}^3}\phi^s_{u_n} u^2_ndx
 +\mu \frac{p(3-2s)-6}{2ps}\int_{\mathbb{R}^3}|u_n|^pdx\Big\}\\
&= \frac{1}{c^2} \Big\{\frac{6s-3}{4s}\int_{\mathbb{R}^3}\phi^s_{u} u^2dx
 +\mu \frac{p(3-2s)-6}{2ps}\int_{\mathbb{R}^3}|u|^pdx\Big\}
< 0,
\end{split}
\end{equation}
for $\mu>\mu^*_2$ large.

With the help  of  the above step 1, step 2, step 3, we can prove $u_n\to u$ 
strongly in $H^s_r(\mathbb{R}^3)$. Let $\mu>\mu^*_2$, set 
$B:=\lim_{n\to\infty} \|(-\Delta)^{s/2}u_n\|^2_2\geq 0 $, by the weak convergence of 
$u_n\rightharpoonup u$ in  $H^s_r(\mathbb{R}^3)$ and  \eqref{5.166}, one obtains
\begin{equation}\label{5.22}
\begin{split}
&(a+bB)\int_{\mathbb{R}^3}(-\Delta)^{s/2}u(-\Delta)^{s/2}\varphi dx
 +\int_{\mathbb{R}^3}\phi^s_{u}u\varphi dx\\
&-\mu \int_{\mathbb{R}^3}|u|^{q-2}u\varphi dx
 -\int_{\mathbb{R}^3}|u|^{2^*_s-2}u\varphi dx \\
&=\lambda \int_{\mathbb{R}^3} u\varphi dx.
\end{split}
\end{equation}
Therefore, from \eqref{5.2122}, \eqref{5.212}, and \eqref{5.22} and 
$u_n\to u$ in $L^{2^*_s}(\mathbb{R}^3)$, it follows that
\begin{align*} %\label{5.23}
&\lim_{n\to\infty}\Big[\Big(a+b\int_{\mathbb{R}^3}|(-\Delta)^{s/2}u_n|^2dx\Big)
 \int_{\mathbb{R}^3}|(-\Delta)^{s/2}u_n|^2 dx
 +\int_{\mathbb{R}^3}\phi^s_{u_n}{u^2_n}dx
 -\lambda \int_{\mathbb{R}^3} {u^2_n} dx\Big]\\
&=\lim_{n\to\infty}\Big [\mu \int_{\mathbb{R}^3}|u_n|^{p}dx
 +\int_{\mathbb{R}^3}|u_n|^{2^*_s} dx\Big]\\
&=\mu \int_{\mathbb{R}^3}|u|^{p}dx+\int_{\mathbb{R}^3}|u|^{2^*_s} dx\\
&=(a+bB)\int_{\mathbb{R}^3}|(-\Delta)^{s/2}u|^2 dx
 +\int_{\mathbb{R}^3}\phi^s_{u}{u^2}dx
 -\lambda \int_{\mathbb{R}^3} {u^2} dx.
\end{align*}
Since $\lambda<0$, as in the proof of Lemma {\ref{lem2.161}},  
we can derive that
\begin{align*} %\label{5.24}
\lim_{n\to\infty} \int_{\mathbb{R}^3} {u^2_n} dx
= \int_{\mathbb{R}^3} {u^2} dx,\quad
\lim_{n\to\infty}  \|(-\Delta)^{s/2}u_n\|^2_2 = \|(-\Delta)^{s/2}u\|^2_2.
\end{align*}
Therefore, $u_n\to u$ in $H^s_{r}(\mathbb{R}^3)$ and $\|u\|_2=c$. 
The proof is complete.
\end{proof}

With the help of the above technical lemmas, we can prove Theorem \ref{thm1.2} 
as follows.

\begin{proof}[Proof of Theorem \ref{thm1.2}]
From Lemmas \ref{lem5.1} and \ref{lem5.2}, the functional $I_\mu$ satisfies 
the Mountain pass geometry. By Lemmas \ref{lem5.4} and \ref{lem5.5}, there exist 
a $(PS)_{c_\mu(c)}$-sequence $\{u_n\}\subset S_{r}(c)$ satisfying  
$P_\mu(u_n)\to 0$ as $n\to \infty$, \eqref{5.16}, and \eqref{5.166}, 
which is bounded in $H^s_r(\mathbb{R}^3)$. 
Furthermore, by Lemma \ref{lem5.7}, there exists $\mu^*_3:=\mu^*_3(c)$ large enough 
such that $0<c_\mu(c)<\Big(\frac{1}{4}-\frac{1}{2^*_s}\Big)(aS)^{\frac{3}{2s}}$ 
for $\mu>\mu^*_3$. Then, by Lemma \ref{rem5.91}, there exist 
$u\in S_r(c)$ and $\lambda<0$ such that passing to a subsequence $u_n\to u$ in 
$H^s_r(\mathbb{R}^3)$ if   $\mu>\mu^*(c):=max\{\mu^*_2, \mu^*_3\}$. 
This completes the proof.
\end{proof}

\subsection*{Acknowledgements}
This research was supported by the Natural Science Foundation of
Sichuan (2023NSFSC0073).
The authors want to express their gratitude to the reviewers for the careful reading and
helpful suggestions which led to an improvement of the original
manuscript.


\begin{thebibliography}{10}

\bibitem{A1111}  Ambrosio,  V.;
\emph{An existence result for a fractional Kirchhoff-Schr\"odinger-Poisson system},
 Z. Angew. Math. Phys., 2018, 69(2),  30.

\bibitem{A5878} Alves, C. O.; Ji, C,; Miyagaki, O. H.;
\emph{Normalized solutions for a Schr\"odinger  equation with critical growth in 
$\mathbb{R}^N$}, Calc. Var. Partial Differential Equations, 2022,   61(1), 18.

\bibitem{P9899}   Bogachev. V. I.;
\emph{Measure Theory}, vol. II, Springer-Verlag, Berlin, 2007.

\bibitem{Bartsch99} Bartsch, T.; De Valeriola. S.;
\emph{Normalized solutions of nonlinear Schr\"odinger equations},
 Arch. Math. (Basel), 2013,  100(1), 75--83.

\bibitem{C4998}   Cho, Y.;  Hwang,  G.;  Kwon, S.;  Lee, S.;
\emph{On finite time blow-up for the mass-critical Hartree equations},
Proc. Roy. Soc. Edinburgh Sect. A,  2015,  14(3), 467--479.

\bibitem{L1159} Che, G. F.; Chen. H. B.; Shi, H. X.; Wang, Z. W.;
\emph{Existence of nontrial solutions for fractional Schr\"odinger-Poisson 
system with sign-changing  potentials}, 
Math. Methods Appl. Sci., 2018, 41(13), 5050--5064.

\bibitem{D7998}  Di Nezza, E.; Palatucci, G.; Valdinoci, E.;
\emph{Hitchhiker's guide to the fractional Sobolev spaces}. 
Bull. Sci. Math., 2012, 136(5), 521-573.

\bibitem{D1678} d'Avenia, P.;  Siciliano, G.;
\emph{Nonlinear Schr\"odinger equation in the Bopp-Podolsky electrodynamics: 
solutions in the electrostatic case}, J. Differ. Equa.,  2019,  267(2), 1025--1065.

\bibitem{D1688} Dinh, V. D.;
\emph{Existence, non-existence and blow-up behaviour of minimizers for the 
mass-critical fractional non-linear Schr\"odinger equations with periodic potentials},
Proc. Roy. Soc. Edinburgh Sect. A,  2020,  150(6), 3252--3292.

\bibitem{D1399}  Duan, X. L.; Wei, G. M.; Yang, H. T.;
\emph{Ground states for a fractional Schr\"odinger-Poisson system involving 
Hardy potentials}, Appl. Anal., 2022, 101(4), 1225--1242.

\bibitem{F1125}  Feng, S. H.; Wang, L.; Huang, L.;
\emph{Least energy sign-changing solutions of fractional Kirchhoff-Schr\"odinger-Poisson 
system with critical and logarithmic nonlinearity}, 
Complex Var. Elliptic Equ., 2023,  68(1), 81--106.

\bibitem{F1114}  Feng, S. H.;  Chen, J. H.; Sun, J. J.; Huang, X. J.;
\emph{Least energy sign-changing solutions for fractional critical 
Kirchhoff-Schr\"odinger-Poisson with steep potential well},
 Fract. Calc. Appl. Anal., 2024, 27(1), 124--156.

\bibitem{G1312}  Guo, L.;
\emph{Sign-changing solutions for fractional Schr\"odinger-Poisson system in
 $\mathbb{R}^3$}, Appl. Anal., 2019, 98(11), 2085--2104.

\bibitem{H2588} He, X. M.; Meng, Y. X.; Squassina, M.;
\emph{Normalized solutions for a fractional Schr\"odinger-Poisson system with 
critical growth}, Calc. Var. Partial Differential Equations, 2024,  63(6), 142.

\bibitem{G1314} Huang, X. Q.; Liao, J. F.;
\emph{Existence and asymptotic behavior for ground state sign-changing solutions 
of fractional Schr\"odinger-Poisson system with steep potential well},
Commun. Anal. Mech., 2024, 16(2), 307--333.

\bibitem{J7888}  Jeanjean, L.;
\emph{Existence of solutions with prescribed norm for semilinear elliptic equations},
Nonlinear Anal., 1997,  28(10), 1633--1659.

\bibitem{L5898} Jeanjean, L.; Lu, S. S.;
\emph{Nonradial normalized solutions for nonlinear scalar field equations},
Nonlinearity, 2019, 32(12), 4942--4966.

\bibitem{J1115}  Jian, H.; Zhong, Q. C.; Wang, L.;
\emph{Positive solutions for fractional Kirchhoff-Schr\"odinger-Poisson system with 
steep potential well},  Rev. Math. Phys.,  2023, 35(07), 2350024.

\bibitem{L7889} Lions, P. L.;
\emph{The concentration-compactness principle in the calculus of variations:
The locally compact case, II},
Ann. Inst. H. Poincar\'e Anal. Non Lin\'eaire, 1984,  1(4), 223--283.

\bibitem{L2998}   Laskin, N.;
\emph{Fractional quantum mechanics and L\'evy path integrals},
Phys. Lett. A, 2000,  268(4--6), 298--305.

\bibitem{L3998}  Laskin, N.;
\emph{Fractional Schr\"odinger equation}, Phys. Rev. E, 2002, 66(5), 56--108.

\bibitem{luo288}  Luo, X.;  Wang, Q. F.;
\emph{Existence and asymptotic behavior of high energy normalized solutions for 
the Kirchhoff type equations in $\mathbb{R}^3$},
 Nonlinear Anal. Real World Appl., 2017, 33, 19--32.

\bibitem{M6998}  Molica B., G.; Radulescu, V. D.; Servadei, R.;
\emph{Variational methods for nonlocal fractional problems. Encyclopedia of 
Mathematics and its Applications}, 162. Cambrige University Press, Cambrige, 2016.

\bibitem{M1113}   Meng, Y. X.; Zhang, X. R.; He, X. M.;
\emph{Least energy sign-changing solutions for a class of fractional Kirchhoff-Poisson 
system}, J. Math. Phys., 2021, 62(9), 091508.

\bibitem{T7899} Teng, K. M.;
\emph{Existence of ground state solutions for the nonlinear fractional 
Schr\"odinger-Poisson system with critical Sobolev exponent},
J. Differential Equations, 2016, 261(6),  3061--3106.

\bibitem{P1199}  Teng, K. M.;
\emph{Ground state solutions for the non-linear fractional Schr\"odinger-Poisson system},
Appl. Anal., 2019, 98(11),  1959--1996.

\bibitem{W6888}  Willem, M.;
\emph{Minimax Theorems}, Birkhauser, Boston 1996.

\bibitem{W1145}  Wang, L.; R\v{a}dulescu, V. D.; Zhang, B. L.;
\emph{Infinitely many solutions for fractional Kirchhoff-Schr\"odinger-Poisson systems},
J. Math. Phys., 2019, 60(1), 011506.

\bibitem{W1135}  Wang, D. B.; Zhang, J. L.;
\emph{Least energy sign-changing solutions of fractional Kirchhoff-Schr\"odinger-Poisson 
system with critical growth}, Appl. Math. Lett., 2020,  106, 106372.

\bibitem{W1112}  Wang, L.;  Han, T.;  Cheng, K.; Wang, J. X.;
\emph{Ground state solutions for nonlinear fractional Kirchhoff-Schr\"odinger-Poisson 
systems}, Int. J. Nonlinear Sci. Numer. Simul., 2021,  22(5), 531--542.

\bibitem{wang1388} Wang, L.; Tang, L. Q.; Wang, J.; Wang, J. X.;
\emph{Existence of normalized solutions for fractional Kirchhoff-Schr\"odinger-Poisson 
equations with general nonlinearities},
 Math. Methods Appl. Sci., 2024, 47(11), 8786--8807.

\bibitem{H2998} Xiang, M. Q.; Wang, F. L.;
\emph{Fractional Schr\"odinger-Poisson-Kirchhoff type systems involving critical 
nonlinearities}, Nonlinear Anal., 2017, 164, 1--26.

\bibitem{yang1288} Yang, J. F.; Guo, W.; Li, W. M.; Zhang, J. F.;
\emph{Existence of normalized solutions for a class of Kirchhoff-Schr\"odinger-Poisson 
equations in $\mathbb{R}^3$}, Ann. Func. Anal., 2023, 14(1), 13.

\bibitem{Z1688} Zhang, X.; Zhang, B. L.; Repov\v{s}, D.;
\emph{Existence and symmetry of solutions for critical fractional Schr\"odinger
 equations with bounded potentials}, Nonlinear Anal., 2016, 142, 48--68.

\end{thebibliography}

\end{document}
